\documentclass[11pt,a4paper]{article}
\usepackage[utf8]{inputenc}
\usepackage[T1]{fontenc}
\usepackage[english]{babel}
\usepackage{amsmath,amsfonts,amssymb}
\usepackage{graphicx}
\usepackage{hyperref}
\usepackage{geometry}
\geometry{margin=2.5cm}
\usepackage{abstract}
\renewcommand{\abstractnamefont}{\normalfont\bfseries}
\usepackage{booktabs}
\usepackage{multirow}
\usepackage{natbib}
\usepackage{parskip}\documentclass[11pt,a4paper]{article}
\usepackage[utf8]{inputenc}
\usepackage[T1]{fontenc}
\usepackage[english]{babel}
\usepackage{amsmath,amsfonts,amssymb}
\usepackage{graphicx}
\usepackage{hyperref}
\usepackage{geometry}
\geometry{margin=2.5cm}
\usepackage{abstract}
\renewcommand{\abstractnamefont}{\normalfont\bfseries}
\usepackage{booktabs}
\usepackage{multirow}
\usepackage{natbib}
\usepackage{parskip}
\usepackage{subcaption}
\usepackage{tikz}
\usepackage{placeins}
\usepackage{float}

% Evitar que figuras/tablas acumuladas pasen al final del documento (p.\ ej.\ tras la bibliografía).
\setcounter{topnumber}{9}
\setcounter{bottomnumber}{9}
\setcounter{totalnumber}{20}
\renewcommand{\topfraction}{0.95}
\renewcommand{\bottomfraction}{0.95}
\renewcommand{\textfraction}{0.07}
\renewcommand{\floatpagefraction}{0.85}

\title{Empirical Mode Decomposition and Graph Transformation of the MSCI World Index: A Multiscale Topological Analysis for Graph Neural Network Modeling}

\author{Agustín Manuel de los Riscos Mayorga}

\date{\today}

\AtBeginDocument{\renewcommand{\abstractname}{}}

\begin{document}

\maketitle

\begin{abstract}
This work targets graph representations that can feed graph neural network (GNN) pipelines and therefore applies Empirical Mode Decomposition (EMD) to the MSCI World index, assigning a graph object to each extracted mode. The exposition begins with controlled synthetic benchmarks, including a multi-component stress test with an equity-like upward drift, where classical EMD is compared with Ensemble EMD (EEMD) and Complete Ensemble EMD with Adaptive Noise (CEEMDAN), using ensemble sizes and noise hyperparameters from a small grid search aimed at low average cross-IMF correlation together with tight numerical reconstruction. That comparison clarifies separation limits, mode mixing, and the trade-offs of noise-assisted schemes and therefore grounds the subsequent choice of CEEMDAN for the empirical series, which yields eight intrinsic mode functions (IMFs) spanning high-frequency structure through slower swings. Once those modes are fixed, three deterministic time-series-to-graph maps (natural visibility (NVG), horizontal visibility (HVG), recurrence graphs) attach a graph instance to each  IMF and, when relevant, to the residue. Because each map stresses different geometry, the induced topology co-varies with both IMF scale and construction family. Visibility graphs emphasize amplitude-driven connectivity and strong local clustering, whereas recurrence graphs track embedding-space proximity and often fragment, so graph-level statistics need not move monotonically with IMF order.

Retaining every oscillatory IMF as its own channel can inflate node and edge budgets under several graph constructions, and a smaller per-date multivariate summary could ease downstream graph workloads or repeated multichannel experiments. Accordingly, the text documents an optional linear compression of the eight oscillatory IMF columns to $k\in\{3,4,5\}$ scores per date via PCA and FastICA on the MSCI CEEMDAN table, with figures and tables focused on $k=4$. That block evaluates reconstructions in native IMF units, reports numerical agreement of ICA and PCA on $\hat{\mathbf{X}}$ at that rank, and contrasts linear coupling among scores with the synthetic cross-IMF Pearson diagnostics, whereas the headline graph statistics in the core results still use native IMFs plus residue without applying that shrinkage. Taking all together, these results show pronounced heterogeneity across scales and encoders, so scale-aware GNN pipelines should adapt to the view at hand rather than reuse one graph template across modes.
\end{abstract}

\vspace{0.5cm}

{\normalfont\normalsize\noindent\textbf{Keywords:} Empirical Mode Decomposition, EEMD, CEEMDAN, MSCI World, Time Series-Graph Transformation, Visibility Graphs, Recurrence Graphs, Graph Neural Networks, Principal Component Analysis, FastICA, Dimensionality Reduction, Financial Time Series}

\newpage
\tableofcontents
\newpage

\section{Introduction}

Analyzing financial time series poses specific challenges derived from the market, whose dynamics are typically neither linear nor stationary. Traditional time series decomposition methods, including wavelet transforms and Fourier analysis, are based on predefined basis functions that often prove inadequate for capturing the inherent properties of financial datasets. The \textbf{Empirical Mode Decomposition} (EMD), originally proposed by \citep{huang1998empirical}, has proven to be highly suitable for the analysis of financial time series, as it provides an adaptive and data-driven methodology that operates without making assumptions about the underlying data structure. Unlike traditional techniques, EMD adaptively locates local extrema and constructs envelopes to isolate oscillatory components, making it highly effective for decomposing financial series. Each of the resulting series is characterized by different time-dependent amplitude and frequency values.

The effectiveness of EMD in the analysis of financial time series stems from this technique's ability to adapt to the fundamental properties that define stock market data: \textbf{non-stationarity} resulting from prolonged trends, \textbf{nonlinear dynamics} evident through volatility clustering and regime transitions, and \textbf{variable amplitude and frequency patterns} according to different market states. These characteristics make financial time series particularly suitable for EMD analysis, given that the technique is explicitly designed to isolate intrinsic mode functions (IMFs) from signals exhibiting these characteristics. The \textbf{MSCI World} index, which serves as a global market indicator, fulfills the properties enumerated above, positioning it as an appropriate case study for applying EMD to the decomposition of time series in this sector.

Although EMD provides a robust decomposition framework, the resulting IMFs require advanced modeling techniques to represent structural and temporal relationships. To achieve this objective, this line of work proposes \textbf{graph neural networks} (GNNs), given their ability to effectively learn from graph-structured data and identify both localized and global patterns through message-passing mechanisms. However, to utilize GNNs in the analysis of IMFs, it is essential to represent temporal sequences as nodes and edges through data transformations, while maintaining the temporal order of the series at all times. Embedding each IMF as its own graph on the full daily grid multiplies structural objects with the number of modes and can inflate node and edge budgets for visibility-type constructions, which motivates documenting and testing optional linear compression of the oscillatory IMF block after the core graph statistics (Sections~\ref{sec:method_dim_red} and~\ref{sec:dim_red_imf}).

The central thread is CEEMDAN on MSCI prices together with complementary graph encodings of each extracted mode. Concretely, the manuscript covers the following points.
\begin{itemize}
    \item \textbf{validating the applicability conditions of EMD} for the MSCI World index through statistical tests.
    \item \textbf{comparing EMD, EEMD, and CEEMDAN on synthetic signals} of controlled design (Section~\ref{sec:synthetic_emd}), with ensemble hyperparameters for the noise-assisted methods tuned on the multi-component family so that the numerical experiment matches the stress test used in the figure.
    \item \textbf{decomposing the index} into IMFs and a residue via CEEMDAN, with a characterization of the energy, variance, frequency, and amplitude attributes of each oscillatory mode, and \textbf{a direct comparison} on the same MSCI closing prices between CEEMDAN and EEMD in terms of the reconstruction gap when the residue is excluded from the IMF sum (Figure~\ref{fig:imf_decomposition}), to highlight the improvement of stage-wise adaptive noise over input-only ensemble averaging for isolating the slow trend.
    \item \textbf{converting each IMF (and the residue, where applicable)} into graph representations using NVG, HVG, and recurrence constructions.
    \item \textbf{calculating metrics of the obtained graphs} (clustering coefficients, centrality measures, connectivity patterns) that summarize the structural characteristics of each decomposed component.
    \item \textbf{documenting optional PCA and FastICA compression} of the eight oscillatory IMF channels (Section~\ref{sec:method_dim_red}), with empirical reconstruction and score-dependence diagnostics on the MSCI CEEMDAN panel after those graph analyses (Section~\ref{sec:dim_red_imf}), motivated by large node and edge budgets when every mode is embedded separately, foregrounding $k=4$ among ranks $k\in\{3,4,5\}$, mapping fits back to native IMF units for RMSE and $R^2$ by mode, and contrasting Pearson coupling among reduced scores with the synthetic cross-IMF summaries, while the tabulated graph constructions use native IMF$_1$--IMF$_8$ and the residue without applying that reduction.
\end{itemize}


Beyond this opening section, Section~\ref{sec:methodology} collects data description, EMD theory and variants, applicability checks, NVG/HVG/recurrence graph rules, and optional PCA and FastICA compression of the oscillatory IMF panel with the graph-motivated rationale developed in Section~\ref{sec:method_dim_red}. Section~\ref{sec:results} orders the empirical blocks as follows. First come suitability tests for EMD on MSCI (Section~\ref{sec:suitcondemd}). Next is the synthetic EMD versus EEMD versus CEEMDAN comparison (Section~\ref{sec:synthetic_emd}). The third block reports CEEMDAN on the index with IMF diagnostics. The fourth summarizes graph statistics under the three maps. The fifth presents the MSCI compression study (Section~\ref{sec:dim_red_imf}) after the graph discussion, tying rank-$k$ linear summaries to the observed node and edge burden, with per-mode reconstruction in native IMF units, auxiliary ranks $k\in\{3,5\}$ beside the $k=4$ exhibits, and a Pearson check on reduced scores parallel to the synthetic IMF panel. Section~\ref{conclfutu} closes with synthesis and forward-looking items.
\section{Methodology}
\label{sec:methodology}

\subsection{Description of the data}

The analysis is based on the daily closing prices of the MSCI World index (Figure~\ref{fig:msci}). The analyzed dataset spans from January 12, 2012 to April 20, 2026, and comprises 3587 daily observations. This period covers various market regimes, including bull markets, corrections, and periods of high volatility. The price series ranges from a minimum of 37.73\$ in US dollars to a maximum of 195.79\$, representing a cumulative return of 407.52 \% during the analysis period. Daily returns are calculated as percentage variations, with a mean daily return of 0.0511 \% and a standard deviation of 1.0726 \%. The 30-day moving average volatility is 0.9380 \%, with a maximum value of 5.1228 \% during periods of high volatility.

\begin{figure}[H]
    \centering
    \begin{tikzpicture}
        \node[anchor=south west,inner sep=0] (image) at (0,0) {\includegraphics[width=\textwidth]{./images/english/msci_world_data.png}};
        \begin{scope}[x={(image.south east)},y={(image.north west)}]
            % Etiqueta a) en el panel superior - esquina superior izquierda
            \node[anchor=north west, font=\bfseries, fill=white, inner sep=2pt, opacity=0.8] at (-0.01,0.98) {(a)};
            % Etiqueta b) en el panel medio - esquina superior izquierda
            \node[anchor=north west, font=\bfseries, fill=white, inner sep=2pt, opacity=0.8] at (-0.01,0.65) {(b)};
            % Etiqueta c) en el panel inferior - esquina superior izquierda
            \node[anchor=north west, font=\bfseries, fill=white, inner sep=2pt, opacity=0.8] at (-0.01,0.32) {(c)};
        \end{scope}
    \end{tikzpicture}
    \caption{Analyzed data of the MSCI World. \textbf{a)} Time series trends of closing prices for the MSCI World index from January 2012 to April 2026, accompanied by the 50 and 200-day moving averages that allow identifying short and long-term patterns. \textbf{b)} Daily returns calculated as percentage variations, which capture the daily variability of the index.  \textbf{c)} 30-day moving volatility, which measures the degree of variation of returns in a 30-day moving window. This metric allows identifying periods of higher risk and market stability, showing pronounced peaks during market events such as those observed in 2020.}
    \label{fig:msci} % Para referenciarla en el texto con \ref{fig:mi_imagen}
\end{figure}
\subsection{Empirical Mode Decomposition. Theory and variants}

The EMD decomposition, introduced by \citep{huang1998empirical}, is an adaptive data-driven method for decomposing non-stationary and nonlinear time series into a finite set of oscillatory components called intrinsic mode functions (IMF). Unlike traditional decomposition methods that rely on predefined basis functions (for example, Fourier transforms that use sinusoidal bases or wavelet transforms that use predefined mother wavelets), EMD adaptively identifies local extrema and constructs envelopes to extract oscillatory components at different temporal scales.

The fundamental principle of EMD is the sifting process, which iteratively extracts IMFs by identifying local maxima and minima, constructing upper and lower envelopes through cubic spline interpolation, calculating the mean envelope, and subtracting it from the signal. This process continues until the resulting signal satisfies the IMF conditions: the number of extrema and zero crossings must differ by at most one, and the mean value of the upper and lower envelopes must be zero at any point. This process ends when a stopping criterion is met, typically based on a standard deviation threshold between consecutive sifting iterations.

However, standard EMD has several limitations, among which the mode mixing problem stands out. When decomposing with the original EMD method, oscillations of different scales tend to mix into a single IMF and similar oscillations may be split across several IMFs. This problem is magnified in the presence of intermittent signals and noise, which can make the decomposition unstable and non-unique.

The \textbf{Complete Ensemble Empirical Mode Decomposition with Adaptive Noise (CEEMDAN)} \citep{torres2011complete} addresses mode mixing by injecting white noise in an adaptive way at each stage of the decomposition. Noise realizations are drawn in the ensemble, the corresponding IMFs are averaged across trials, and the procedure continues until the full decomposition—including a residual trend—is obtained. Compared with injecting noise only at the first stage, this construction typically yields more stable IMFs with lower residual noise contamination for a given ensemble size.

To obtain good results in the decomposition of the MSCI World, CEEMDAN is applied with the execution parameters defined in Table~\ref{tab:ceemdan_params}. Those values are specific to the empirical series. The exploratory synthetic experiments in Section~\ref{sec:synthetic_emd} use the ensemble sizes and noise hyperparameters reported there (including a grid search on the multi-component stress test), which need not coincide numerically with Table~\ref{tab:ceemdan_params}.

\begin{table}[h]
\centering
\caption{CEEMDAN decomposition parameters}
\label{tab:ceemdan_params}
\begin{tabular}{ll}
\toprule
\textbf{Parameter} & \textbf{Value} \\
\midrule
Maximum number of IMFs & 14 \\
Ensemble trials & 100 \\
Noise amplitude $\varepsilon$ & 0.05 (5\% of signal's std dev) \\
Pseudorandom seed (noise) & 42 \\
\bottomrule
\end{tabular}
\end{table}

\subsection{Prerequisites for effective EMD application.}

For EMD analysis to be effective, the time series must satisfy the following key conditions:

\begin{itemize}
\item \textbf{Non-stationarity}: EMD decomposition offers optimal performance with non-stationary series, as it can adaptively capture trends and components at different temporal scales. Therefore, the stationarity of the studied MSCI World time interval is checked using the Augmented Dickey-Fuller (ADF) test \citep{dickey1979distribution} and the Kwiatkowski-Phillips-Schmidt-Shin (KPSS) test \citep{kwiatkowski1992testing}.

\item \textbf{Non-linearity}: Applying EMD techniques to the series does not require linearity assumptions. Non-linearity is evaluated through multiple approaches: using the Brock-Dechert-Scheinkman (BDS) test \citep{brock1996test} on the residuals of an AR(1) model, through autocorrelation analysis of squared returns to detect ARCH/GARCH effects \citep{engle1982autoregressive,bollerslev1986generalized}, and finally using the Ljung-Box test \citep{ljung1978measure} on the squared returns of the index.

\item \textbf{Variable amplitude and frequency}: EMD excels at capturing modes with time-varying amplitude and frequency. Therefore, amplitude variability is analyzed through continuous volatility measures (252-day windows corresponding to one trading year) and frequency variability through zero-crossing analysis.
\end{itemize}

The results of all tests listed in this section are described in Section~\ref{sec:suitcondemd}.

\subsection{Time series to graph transformations}

The transformation of time series into graph representations allows us to analyze the structural properties of temporal data using graph theory techniques and network analysis. Furthermore, converting to graphs enables the application of graph neural networks and offers an alternative perspective for understanding the underlying dynamics of decomposed time series components. In this study, we employ three complementary graph transformation methods to convert the extracted IMFs into graph structures: the natural visibility graph (NVG), the horizontal visibility graph (HVG), and recurrence-based networks.

\subsubsection{Visibility algorithms}

Visibility graph algorithms, introduced by \citep{lacasa2008visibility}, transform a time series into a graph by establishing connections between data points based on geometric visibility criteria. The fundamental principle is that two points in the time series are connected if no intermediate point obstructs the line of sight between them.

The variant called the natural visibility graph (NVG) algorithm constructs a graph in which each data point $(t_i, y_i)$ of the time series becomes a node, and an edge exists between nodes $i$ and $j$ if, for all intermediate points $k$ such that $i < k < j$, the following condition is satisfied:
\begin{equation}
y_k < y_i + (y_j - y_i) \frac{t_k - t_i}{t_j - t_i}
\end{equation}

This geometric criterion ensures that the line segment connecting points $i$ and $j$ lies above all intermediate data points, establishing a natural visibility relationship. NVG preserves important properties of the original time series, such as periodicity, fractality, and hierarchical structure, making it suitable for analyzing complex temporal patterns in financial data.

In a second variant, the horizontal visibility graph (HVG) algorithm \citep{luque2009horizontal} simplifies NVG by imposing a stricter horizontal visibility constraint. Two nodes $i$ and $j$ are connected if, for all intermediate points $k$ with $i < k < j$, the following condition is satisfied:
\begin{equation}
y_k < \min(y_i, y_j)
\end{equation}

These two transformations have demonstrated utility in various recent applications in the financial domain. An example is found in Han et al. (2022), where they developed MOHVGCA (Multiscale Online-Horizontal-Visibility-Graph Correlation Analysis), a method that employs HVG to quantify scale-dependent correlations between financial time series, revealing dynamic associations between WTI oil futures and various currencies during oil crises \citep{han2022multiscale}. More recently, Zeng and Chen (2025) proposed the use of rising visibility graphs to extract topological motifs from trend structures in stock indices, demonstrating that these recurrent topological structures can be identified and used to forecast emerging patterns in financial time series \citep{chen2025identifying}. These studies demonstrate the ability of visibility algorithms to reveal hidden structural characteristics in time series from various domains, validating their application in the analysis of complex financial data as addressed in this work.

\subsubsection{Recurrence-based algorithms}

As an alternative to visibility algorithms, the recurrence matrix can be employed as an adjacency matrix to transform time series into graph representations by identifying recurrence patterns within the dynamical system \citep{eckmann1987recurrence}. This method captures the recurrence structure of the series in a higher-dimensional embedding space. This effectively reveals both the geometry of the underlying attractor and the essential dynamical properties of the system.

The construction of recurrence graphs begins with the determination of optimal embedding parameters using geometric and information-theoretic criteria. The delay parameter $\tau$ is selected by analyzing mutual information, specifically identifying the first minimum of the mutual information function between the original series and its time-delayed version. This ensures that the delayed coordinates provide maximum information about the system state while minimizing redundancy. Subsequently, the optimal embedding dimension $d$ is estimated using the false nearest neighbors (FNN) method. This technique identifies the minimum dimension required to unfold the attractor by detecting false neighbors that appear close in low dimensions but are separated in higher dimensions. The FNN method systematically increases the embedding dimension until the proportion of false neighbors falls below a threshold, typically 10\%.

Once $\tau$ (time delay) and $d$ (embedding dimension) are determined, the embedding space is constructed using delay coordinate embedding. Each point $X_i$ in this space is defined as a $d$-dimensional vector:
\begin{equation}
X_i = [Y_i, Y_{i+\tau}, Y_{i+2\tau}, \ldots, Y_{i+(d-1)\tau}]
\end{equation}
where $X_i$ represents the $i$-th point in the embedding space and $Y_i$ denotes the value of the original time series at time $i$.

The resulting recurrence graph preserves the topological and geometric properties of the underlying attractor, enabling the analysis of dynamical invariants such as correlation dimension, Lyapunov exponents, and entropy measures. Node features in the recurrence graph are assigned using the first component of the embedding vectors, maintaining a connection with the original time series values.

From the generated embedding space, the recurrence graph is constructed by computing pairwise distances between all embedding vectors. Two nodes $i$ and $j$ in the graph are connected if the distance between their corresponding embedding vectors $X_i$ and $X_j$ is below a threshold $\varepsilon$. This is formally expressed through the $(i, j)$ element of the recurrence matrix $R$:
\begin{equation}
R_{ij} = \begin{cases}
1 & \text{if } \left\|X_i - X_j\right\| \le \varepsilon \\
0 & \text{otherwise}
\end{cases}
\end{equation}
where $R_{ij}$ is the $(i, j)$-th element of the recurrence matrix. The threshold $\varepsilon$ is typically chosen as a percentile of the pairwise distance distribution, ensuring that the graph captures the most significant recurrent structures while maintaining computational tractability.

Numerous studies have demonstrated that the transformation of time series to recurrence graphs has proven particularly useful in financial market analysis, where capturing nonlinear dynamics and complex patterns is crucial for prediction and anomaly detection. In this context, a hybrid dual-branch architecture combines recurrence plots with transposed transformers for stock trend prediction \citep{su2025hybrid}. In that setup, closing price time series are transformed into recurrence plots using sliding windows of historical segments. The first branch generates RP images from these univariate sequences and extracts recurrence quantification analysis (RQA) measures, which capture nonlinear characteristics and temporal relationships between points. These RQA measures include metrics such as recurrence rate (RR), determinism (DET), and laminarity (LAM), which quantify the predictability and structure of the dynamical system. The second branch processes multivariate time series features using a transposed transformer, and both branches are merged to classify the next-day trend, either up or down. Evaluation on seven selected stocks indicates that combining recurrence-based features with sequence modeling through transformers can yield higher accuracy and F1-scores than traditional machine learning and deep learning baselines. This line of work illustrates how recurrence plots can capture nonlinear dynamics and relationships between temporal points that are difficult to detect using conventional time series analysis methods.

Separately, multiplex recurrence networks (MRNs) have been used to search for early warning signals of critical transitions and crashes in stock markets \citep{song2024early}. In that framework, multidimensional return vectors are transformed into multiplex recurrence networks, where each dimension, corresponding to different assets or indicators, forms a network layer. Recurrence relationships are established within each layer and between layers, capturing both the individual dynamics of each asset and the interdependencies between them. Topological summaries of MRNs, specifically average mutual information and average edge overlap, track changes in the recurrence structure over time. Those summaries are assessed on simulated systems and on empirical data from Chinese and US markets, where indicators derived from MRNs appear to anticipate imminent state transitions and to shed light on major historical market episodes.

Following the review of transformation types, it is concluded that applying visibility and recurrence algorithms to IMF (Intrinsic Mode Function) components is of considerable interest. This methodological strategy enables a more comprehensive analysis of the resulting graph, as visibility complements information regarding the temporal structure of the sequence, while recurrence reveals the relationships of the dynamical system inherent to the time series. As discussed in Section~\ref{conclfutu}, applying other types of transformations will be considered in future research lines.

\subsection{Linear dimensionality reduction of a multichannel IMF panel}
\label{sec:method_dim_red}

An EMD-type decomposition turns one observed series into several aligned oscillatory channels indexed by mode order. That expansion is valuable for scale-resolved diagnostics, but in the graph-oriented workflow emphasized in this article it also multiplies the objects passed to later stages: each IMF can be embedded as its own graph, and aggregate cost then scales with the number of modes as well as with the node and edge burden of the chosen visibility or recurrence construction (Section~\ref{sec:results}). Linear reduction of the IMF panel is treated as an \emph{optional} companion: usefulness is checked empirically rather than postulated. The target is to see whether a handful of derived coordinates per date can carry enough joint IMF variation when multivariate width, repeated experiment cost, or encoder size must be capped, at the cost of losing one-to-one alignment with individual IMFs.

Let $\mathbf{X}\in\mathbb{R}^{T\times p}$ denote a panel of $p$ aligned oscillatory IMFs (columns) over $T$ time indices (rows), obtained after an EMD-type decomposition technique. The aim is to obtain, at each time row and without temporal subsampling, a smaller set of $k<p$ derived coordinates by applying a linear map from $\mathbb{R}^{p}$ to $\mathbb{R}^{k}$ rowwise after a fixed affine preprocessing of the columns. Because this construction linearly recombines IMF amplitudes, it sacrifices the scale-specific interpretation of individual IMFs under EMD. When graph construction or repeated experiments constrain multichannel input size, however, it offers a compact multivariate summary of the oscillatory block. To keep the slow trend from contaminating the reduced coordinates, any EMD residue is excluded from the linear block whenever it is present. To place IMFs on a comparable scale before the rowwise map, the oscillatory columns are standardized (centering and unit-variance scaling) prior to projection, following the affine convention used in the PCA and ICA subsections below.

\subsubsection{Principal component analysis (PCA)}
PCA seeks orthogonal directions of maximal variance after centering the columns (here, after the standardization step above) \citep{jolliffe2016principal}. The resulting score matrix $\mathbf{Z}\in\mathbb{R}^{T\times k}$ has columns that are pairwise uncorrelated by construction and orders directions by explained variance in the standardized feature space. Retaining only $k<p$ components is the best rank-$k$ approximation to the standardized data in the Frobenius norm (equivalently, in the least-squares sense): an orthogonal projection onto a $k$-dimensional subspace. Denoting the reconstructed oscillatory block by $\hat{\mathbf{X}}$, the approximation error $\mathbf{X}-\hat{\mathbf{X}}$ is generally nonzero whenever $k<p$.

\subsubsection{Independent component analysis (ICA)}
ICA postulates a linear mixing model in which each row of the observed (whitened) matrix is a linear combination of latent sources that are mutually \emph{statistically independent} rather than merely uncorrelated \citep{hyvarinen2000independent}. Estimation typically proceeds after whitening (decorrelation and unit variance) and uses measures of non-Gaussianity or contrast functions related to mutual information. FastICA is a standard fixed-point algorithm in this family \citep{hyvarinen2000independent}. When the mixing is invertible and the sources satisfy the usual identifiability conditions, ICA links to blind source separation. In the present setting it serves as a complementary linear reduction whose scores are driven toward independence. For a given $k$, the induced reconstruction of the standardized IMF block need not coincide with that of PCA, although both are linear projections of the same column-standardized panel.

\section{Results}
\label{sec:results}

\subsection{Suitable conditions for EMD application.}
\label{sec:suitcondemd}

A battery of tests is performed to confirm an adequate study via EMD decomposition of the indexed series. Our comprehensive evaluation of the MSCI World index confirms that the time series meets the three fundamental conditions for effective EMD analysis. The results of the three evaluated conditions are presented below.

\subsubsection{Stationarity}
Stationarity tests provide strong evidence of the non-stationarity of the price series, which is ideal for EMD. The Augmented Dickey-Fuller (ADF) test on the closing price series yields a test statistic of 1.6628 with a p-value of 0.9980, which does not allow rejecting the null hypothesis of a unit root at all conventional significance levels. The critical values at the 1\%, 5\% and 10\% levels are -3.4322, -2.8624 and -2.5672, respectively, all substantially more negative than the test statistic. The Kwiatkowski-Phillips-Schmidt-Shin (KPSS) test, which has stationarity as the null hypothesis, produces a test statistic of 8.3015 with a p-value of 0.0100, which rejects the null hypothesis of stationarity at the 5\% significance level. Conversely, daily returns are confirmed as stationary in both tests: ADF statistic of -19.7278 (p-value $<$ 0.0001) and KPSS statistic of 0.0447 (p-value = 0.1000). These complementary results confirm that the price series exhibits non-stationary behavior driven by long-term trends, making it highly suitable for EMD decomposition.

\subsubsection{Linearity}
Multiple tests provide consistent evidence of strong nonlinear dynamics in the MSCI World index. The Brock-Dechert-Scheinkman (BDS) test, applied to the residuals of an AR(1) model, detects significant nonlinear dependencies across multiple embedding dimensions. In dimensions 3, 4 and 5, the BDS test statistics are 20.99, 24.95 and 27.83, respectively, all with p-values $<$ 0.0001, which strongly rejects the null hypothesis of independence and indicates the presence of a nonlinear structure.

The distributional analysis further supports the nonlinear characterization. The price series shows significant deviations from normality, with a skewness of 0.8204 and a kurtosis of -0.2243. The Jarque-Bera test statistic yields a p-value $<$ 0.0001, which strongly rejects the null hypothesis of normality. These findings collectively demonstrate that the MSCI World index exhibits rich nonlinear dynamics, making EMD an adequate decomposition method.

\subsubsection{Volatility in amplitude and frequency}

The autocorrelation analysis of squared returns reveals substantial volatility clustering, a distinctive feature of nonlinear ARCH/GARCH effects. We observe 20 significant autocorrelation lags out of the 20 analyzed, with a maximum autocorrelation of 0.9978, indicating strong volatility persistence. The Ljung-Box test on squared returns produces a large test statistic with a p-value $<$ 0.0001, which provides overwhelming evidence of ARCH/GARCH effects and confirms nonlinear dependence in the variance process.

\subsection{Exploratory comparison of EMD-type methods on synthetic signals}
\label{sec:synthetic_emd}

Before applying decomposition techniques to the MSCI World index, we compared three decomposition families---\textbf{EMD}, \textbf{EEMD} \citep{wu2009ensemble}, and \textbf{CEEMDAN} \citep{torres2011complete}---on \emph{synthetic} time series with known structure. The goal is to stress stylized regimes linked to mode mixing and scale separation: closely spaced harmonics (where amplitude beating stresses frequency-resolution limits \citep{rilling2008one}), linearly swept-frequency chirps (non-stationary instantaneous frequency), intermittent high-frequency bursts on a low-frequency carrier, a compact superposition with additive Gaussian noise, and---as an aggregate stress test---a \emph{multi-component} waveform that sums mode-mixing, ascending and descending chirps, close tones, two sinusoids at $0.5\,\mathrm{Hz}$ and $3\,\mathrm{Hz}$, stronger Gaussian noise, and a slow monotonic uptrend (linear ramp) to mimic the cumulative drift of an equity index. This exercise is illustrative rather than a formal identifiability study. It informs how aggressively each method allocates energy across IMFs when the ground-truth components are not orthogonal harmonics. Figure~\ref{fig:synthetic_signals} shows the five synthetic families (duration $5\,\mathrm{s}$, sampling rate $500\,\mathrm{Hz}$), generated with the same parameter settings as the numerical experiments reported below. The exposition first compares native IMF panels (Table~\ref{tab:synthetic_emd_summary}), then applies FastICA to the same stored decompositions (Table~\ref{tab:synthetic_ica_summary}), and finally states how both blocks guide the CEEMDAN+ICA choice on the empirical index.

\begin{figure}[H]
    \centering
    \includegraphics[width=\textwidth]{./images/english/synthetic_signals_emdsynth.png}
    \caption{Synthetic signals used for the exploratory comparison of EMD-type methods. \textbf{(a)} Sum of two sinusoids with nearby frequencies ($10\,\mathrm{Hz}$ and $11.3\,\mathrm{Hz}$), producing amplitude beating. \textbf{(b)} Linear chirp with initial frequency $1.5\,\mathrm{Hz}$ and sweep rate $1.2\,\mathrm{Hz\,s^{-1}}$. \textbf{(c)} Low-frequency carrier ($0.8\,\mathrm{Hz}$) plus a higher-frequency burst ($12\,\mathrm{Hz}$) active only on $t\in[1.2, 3.2]\,\mathrm{s}$. \textbf{(d)} Superposition of the close-tone signal, $0.35$ times the chirp from panel (b), and additive Gaussian noise scaled to $3\%$ of the standard deviation of the close-tone component. \textbf{(e)} Multi-component benchmark that aggregates the elementary generators: mode mixing ($0.3$--$7\,\mathrm{Hz}$ burst on $t\in[1.25, 3.75]\,\mathrm{s}$), ascending chirp ($f_0=0.8\,\mathrm{Hz}$, $k=0.2\,\mathrm{Hz\,s^{-1}}$), descending chirp ($f_0=4\,\mathrm{Hz}$, $k=-0.15\,\mathrm{Hz\,s^{-1}}$), close tones ($1.5$ and $1.7\,\mathrm{Hz}$), sinusoids at $0.5\,\mathrm{Hz}$ and $3\,\mathrm{Hz}$, additive Gaussian noise with standard deviation $0.2$, and a slow linear uptrend from $0$ to $8$ amplitude units over the $5\,\mathrm{s}$ window.}
    \label{fig:synthetic_signals}
\end{figure}

For each family, the same signal was decomposed with EMD, with EEMD (noise-assisted ensemble averaging at the input), and with CEEMDAN (noise-assisted, stage-wise construction aimed at a \emph{complete} ensemble decomposition \citep{torres2011complete}). Unless stated otherwise, we used a shared cap of $K=12$ intrinsic modes (plus the residue as the final component when applicable). Ensemble sizes were set to $M_{\mathrm{EEMD}}=100$ and $M_{\mathrm{CEEMDAN}}=180$, with CEEMDAN noise scale $\varepsilon=0.03$ and EEMD hyperparameters $(\mathrm{noise\_width},\mathrm{SD\_thresh})=(0.095,0.18)$, values selected by a small grid search on the multi-component benchmark (panel (e)) to reduce average cross-IMF correlation while maintaining numerical reconstruction. For EEMD, the residual supplied with the ensemble mean was included so that the IMFs plus the residual reproduce the input up to numerical roundoff in every case. Empirically, relative reconstruction errors $\|\hat{x}-x\|/\|x\|$ stayed below $10^{-15}$ across all runs.

Table~\ref{tab:synthetic_emd_summary} summarizes linear dependence among \emph{native} intrinsic modes. Let $r_{ij}$ denote the sample Pearson correlation between IMF$_i$ and IMF$_j$ over the $T=2500$ samples, for all pairs among the first $\min\{8,N_{\mathrm{IMF}}\}$ intrinsic modes (excluding the residue). For each pair we test the two-sided hypothesis $H_0:\rho_{ij}=0$ using the standard Pearson correlation test (equivalent to a $t$-statistic for $r_{ij}$ with $T-2$ degrees of freedom under classical Gaussian assumptions). We report $\bar{\rho}=\mathrm{mean}_{i<j}|r_{ij}|$ and the proportion $\pi_{0.05}$ of pairs with $p<0.05$ \emph{without} multiple-comparison adjustment across pairs. With large $T$, even modest $|r_{ij}|$ can yield very small $p$-values (high power), so $\pi_{0.05}$ is interpreted jointly with $\bar{\rho}$ rather than as standalone evidence of ``significance'' in the econometric sense.

\begin{table}[ht!]
\centering
\scriptsize
\setlength{\tabcolsep}{3pt}
\caption{Synthetic experiments: $N$ is the number of returned components (including the residue). $\bar{\rho}$ is the mean absolute Pearson correlation between IMF pairs, and $\pi_{0.05}$ is the fraction of pairs with $p<0.05$ (uncorrected for multiple tests). The number of IMF pairs entering these summaries ranges from $3$ to $28$ depending on $N$ and the eight-mode cap. Same settings as Figure~\ref{fig:synthetic_signals}. The block below the double rule isolates the multi-component stress-test (panel (e)). \textbf{Boldface} in that row marks the method with the smallest $\bar{\rho}$ for that benchmark (weakest average cross-IMF linear coupling).}
\label{tab:synthetic_emd_summary}
\begin{tabular}{l*{9}{c}}
\toprule
 & \multicolumn{3}{c}{\textbf{EMD}} & \multicolumn{3}{c}{\textbf{EEMD}} & \multicolumn{3}{c}{\textbf{CEEMDAN}} \\
\cmidrule(lr){2-4} \cmidrule(lr){5-7} \cmidrule(lr){8-10}
\textbf{Family} & $N$ & $\bar{\rho}$ & $\pi_{0.05}$ & $N$ & $\bar{\rho}$ & $\pi_{0.05}$ & $N$ & $\bar{\rho}$ & $\pi_{0.05}$ \\
\midrule
Close tones & 6 & 0.099 & 0.40 & 12 & 0.072 & 0.21 & 8 & 0.090 & 0.57 \\
Linear chirp & 4 & 0.080 & 0.33 & 12 & 0.089 & 0.29 & 7 & 0.071 & 0.47 \\
Burst on carrier & 4 & 0.071 & 0.67 & 11 & 0.103 & 0.29 & 8 & 0.108 & 0.29 \\
Close tones + chirp + noise & 8 & 0.064 & 0.43 & 12 & 0.092 & 0.29 & 8 & 0.087 & 0.19 \\
\addlinespace[3pt]
\midrule
\midrule
\multicolumn{10}{l}{\scriptsize\emph{Stress-test: highly aggregated multi-component signal (Figure~\ref{fig:synthetic_signals}(e))}} \\[2pt]
\textbf{Multi-comp.\ (panel (e))} & 8 & 0.040 & 0.19 & 11 & 0.071 & 0.32 & \textbf{8} & \textbf{0.038} & \textbf{0.24} \\
\bottomrule
\end{tabular}
\end{table}

Because native IMF columns can retain substantial linear dependence even after a noise-assisted decomposition (Table~\ref{tab:synthetic_emd_summary}), we next apply \textbf{FastICA} \citep{hyvarinen2000independent} to the same stored IMF tables. The oscillatory block uses the leading $p=\min\{8,N_{\mathrm{IMF}}\}$ intrinsic modes in IMF index order (matching the cross-IMF window of Table~\ref{tab:synthetic_emd_summary}); the residue is excluded from the fit. FastICA is run with rank $k=\min\{4,p\}$, paralleling the $k=4$ reduction used later on the MSCI CEEMDAN panel (when $p<4$, $k=p$). Column-standardization precedes ICA, as in Section~\ref{sec:method_dim_red}. Table~\ref{tab:synthetic_ica_summary} reports the same pairwise diagnostics on the $k$ estimated sources; the column $N$ lists the post-reduction channel count $k$ plus the unmodified residue when present. The numerical entries use FastICA with pseudorandom seed $42$, sample length $T=2500$, and output dimension $k=\min\{4,p\}$ as above.

\begin{table}[ht!]
\centering
\scriptsize
\setlength{\tabcolsep}{3pt}
\caption{Post-\textbf{FastICA} \emph{with rank reduction} on the synthetic IMF panels from Table~\ref{tab:synthetic_emd_summary}. The estimator reads $p=\min\{8,N_{\mathrm{IMF}}\}$ IMF columns and outputs $k=\min\{4,p\}$ sources (same $k=4$ target as the MSCI ICA block whenever $p\ge 4$). \textbf{$N$} is the channel count \emph{after} reduction: $N=k+1$ when the residue is appended unchanged, otherwise $N=k$. $\bar{\rho}$ and $\pi_{0.05}$ are computed from Pearson correlations among the $k$ ICA sources only (uncorrected pairs, as in Table~\ref{tab:synthetic_emd_summary}). Table entries follow the protocol above (seed $42$, $T=2500$, $k=\min\{4,p\}$); $\bar{\rho}$ and $\pi_{0.05}$ are shown in scientific notation with mantissas rounded to two decimals.}
\label{tab:synthetic_ica_summary}
\begin{tabular}{l*{9}{c}}
\toprule
 & \multicolumn{3}{c}{\textbf{EMD}} & \multicolumn{3}{c}{\textbf{EEMD}} & \multicolumn{3}{c}{\textbf{CEEMDAN}} \\
\cmidrule(lr){2-4} \cmidrule(lr){5-7} \cmidrule(lr){8-10}
\textbf{Family} & $N$ & $\bar{\rho}$ & $\pi_{0.05}$ & $N$ & $\bar{\rho}$ & $\pi_{0.05}$ & $N$ & $\bar{\rho}$ & $\pi_{0.05}$ \\
\midrule
Close tones & 5 & $4.79\times 10^{-16}$ & $0$ & 5 & $8.39\times 10^{-16}$ & $0$ & 5 & $5.61\times 10^{-16}$ & $0$ \\
Linear chirp & 4 & $3.67\times 10^{-16}$ & $0$ & 5 & $1.14\times 10^{-14}$ & $0$ & 5 & $1.06\times 10^{-15}$ & $0$ \\
Burst on carrier & 4 & $3.33\times 10^{-16}$ & $0$ & 5 & $1.07\times 10^{-15}$ & $0$ & 5 & $5.64\times 10^{-16}$ & $0$ \\
Close tones + chirp + noise & 5 & $2.97\times 10^{-15}$ & $0$ & 5 & $9.11\times 10^{-14}$ & $0$ & 5 & $6.06\times 10^{-16}$ & $0$ \\
\addlinespace[3pt]
\midrule
\midrule
\multicolumn{10}{l}{\scriptsize\emph{Stress-test: highly aggregated multi-component signal (Figure~\ref{fig:synthetic_signals}(e))}} \\[2pt]
\textbf{Multi-comp.\ (panel (e))} & 5 & $3.58\times 10^{-15}$ & $0$ & 5 & $4.76\times 10^{-16}$ & $0$ & 5 & $6.34\times 10^{-16}$ & $0$ \\
\bottomrule
\end{tabular}
\end{table}

The native-mode statistics are consistent with the intended roles of the three algorithms. Classical EMD returned the smallest $N$ on the linear chirp and burst-on-carrier families ($N=4$ in both cases), but that parsimony aligns with the known mode-mixing risk under intermittency. On the burst family, the variance share of IMF$_1$ was about $0.37$ for EMD versus order $10^{-3}$ for EEMD and order $10^{-4}$ for CEEMDAN, whereas $\bar{\rho}$ ranked highest for CEEMDAN ($0.108$), then EEMD ($0.103$), then EMD ($0.071$). EEMD often hit the $N=12$ cap on simpler families and stopped at $N=11$ on burst and multi-component rows, while CEEMDAN returned $N\in\{7,8\}$ across the five rows (Table~\ref{tab:synthetic_emd_summary}). On the multi-component benchmark (Figure~\ref{fig:synthetic_signals}(e)), with hyperparameters tuned on that panel, CEEMDAN attained the smallest cross-IMF $\bar{\rho}=0.038$ with $N=8$, slightly below EMD ($0.040$), whereas EEMD remained the most redundant at $\bar{\rho}=0.071$ with $N=11$. Residue variance shares on that stress test were about $0.76$ (CEEMDAN), $0.73$ (EMD), and $0.15$ (EEMD), illustrating how the explicit drift is partitioned across IMFs and the residue.

FastICA on the synthetic panels at rank $k=\min\{4,p\}$ on the same $p=\min\{8,N_{\mathrm{IMF}}\}$ IMF window (Table~\ref{tab:synthetic_ica_summary}) drives mean absolute Pearson correlation between the $k$ estimated sources to numerical roundoff and sets $\pi_{0.05}=0$ for every family and every decomposition backend, so the linear coupling visible in Table~\ref{tab:synthetic_emd_summary} is removed in the reduced subspace whether the extractor was EMD, EEMD, or CEEMDAN. For the \emph{empirical} MSCI World pipeline, the controlled evidence therefore motivates \textbf{CEEMDAN} as the decomposition chosen on the drift-rich stress test, together with \textbf{FastICA on the oscillatory IMF block} when downstream stages require both a compact multivariate summary (the table's $N=k+1$ counts $k$ scores plus the residue, in line with rank-$k$ scores when $k<p$ on MSCI) and channels whose \emph{sample} Pearson correlations between rotated components are effectively zero---not merely a parsimonious panel, but a post hoc cancellation of linear redundancy among decomposed modes that CEEMDAN alone does not fully eliminate. Full statistical independence is stronger than zero linear correlation; the ICA table documents the latter at machine precision. IMF indices should still not be read as uniquely identifying independent economic shocks when the data-generating process is unknown.

\FloatBarrier

\subsection{CEEMDAN decomposition of the MSCI World index}
\label{sec:ceemdan_msci}

The decomposition of the MSCI World index was performed using CEEMDAN with the parameters specified in Table~\ref{tab:ceemdan_params}. This choice targets stable mode separation under noise-assisted ensemble averaging, with adaptive noise injection across decomposition stages \citep{torres2011complete}.

The application of CEEMDAN to the closing price series of the MSCI World index, which comprises 3587 daily observations from January 12, 2012 to April 20, 2026, produced a decomposition into \textbf{eight intrinsic mode functions (IMFs)} and a \textbf{residue}. The residue captures the slow-varying level of the index. It is not required to be strictly monotonic in this implementation, but it carries the bulk of the low-frequency energy removed from the oscillatory IMFs.

\subsubsection{Characterization of the IMFs}

Each extracted IMF represents oscillations at different temporal scales, from high-frequency components (IMF$_1$) to lower-frequency components (IMF$_8$). Table~\ref{tab:imf_metrics} summarizes energy, variance, a simple oscillation-count proxy based on local maxima (reported as ``cycles''), and the mean absolute value as an amplitude summary.

\begin{table}[h]
\centering
\caption{Main metrics of the IMFs extracted via CEEMDAN}
\label{tab:imf_metrics}
\small
\begin{tabular}{lcccc}
\toprule
\textbf{IMF} & \textbf{Energy} & \textbf{Variance} & \textbf{Frequency} & \textbf{Amplitude} \\
 & & & \textbf{(cycles)} & \textbf{Average} \\
\midrule
IMF$_1$ & $7.34 \times 10^2$ & 0.2045 & 1287.0 & 0.3164 \\
IMF$_2$ & $5.14 \times 10^2$ & 0.1433 & 1075.0 & 0.2728 \\
IMF$_3$ & $1.19 \times 10^3$ & 0.3307 & 332.0 & 0.3769 \\
IMF$_4$ & $2.77 \times 10^3$ & 0.7724 & 151.0 & 0.5758 \\
IMF$_5$ & $6.19 \times 10^3$ & 1.7223 & 69.0 & 0.8446 \\
IMF$_6$ & $8.87 \times 10^3$ & 2.4622 & 34.0 & 1.0746 \\
IMF$_7$ & $3.82 \times 10^4$ & 10.5136 & 13.0 & 2.2015 \\
IMF$_8$ & $2.44 \times 10^5$ & 68.0636 & 4.0 & 6.0934 \\
\bottomrule
\end{tabular}
\end{table}

The analysis of the metrics reveals clear patterns in the structure of the IMFs. Lower-order IMFs (IMF$_1$ to IMF$_4$) exhibit high dominant frequencies (from 151 to 1287 cycles) and relatively small amplitudes (all below 0.58), indicating that they capture short-term oscillations and high-frequency noise. As the order of the IMF increases, a systematic decrease in dominant frequency and a progressive increase in variance and average amplitude are observed. Higher-order IMFs (IMF$_7$ and IMF$_8$) present dominant frequencies of 13 and 4 cycles and substantially larger average amplitudes (2.20 and 6.09), reflecting slower, larger-amplitude oscillations at the end of the intrinsic-mode ladder. Among the eight IMFs, IMF$_8$ attains the highest energy ($2.44 \times 10^5$) and variance (68.06), while most of the total signal energy is carried by the residue (not listed in Table~\ref{tab:imf_metrics}), consistent with interpreting the residue as the slow-moving price trend rather than as a fast oscillatory IMF.

\subsubsection{Validation of the decomposition}

Reconstruction was validated by comparing the sum of all IMFs and the residue with the original MSCI World closing-price series. In a consistent decomposition, the original MSCI Index and fully reconstructed signal must be identical.

The CEEMDAN residual is defined as the original signal minus the sum of all eight IMFs. For this sample it has mean 89.93, standard deviation 36.76, and range from 46.20 to 159.99. Figure~\ref{fig:imf_decomposition} (lower panel) contrasts this quantity with the analogous EEMD reconstruction gap on the same prices: the CEEMDAN trace is visually an almost noise-free, monotonically increasing trend, whereas EEMD retains a jagged slow component. That pattern is consistent with attributing the equity-like drift to a single smooth mode under CEEMDAN rather than folding fast oscillations into the baseline. The side-by-side comparison on the MSCI World series thus provides empirical evidence that CEEMDAN yields better decomposition outcomes than EEMD for this index, in the sense of a cleaner separation of slow drift from oscillatory structure on the same closing-price record. Because the CEEMDAN residual behaves here as an almost smooth, monotone trend, it can be modeled to a useful approximation by simple linear regression (for example, regressing the residual on a deterministic time or observation index) when a parsimonious parametric summary of the slow level is needed.

\begin{figure}[H]
    \centering
    \begin{tikzpicture}
        \node[anchor=south west,inner sep=0] (image) at (0,0) {\includegraphics[width=\textwidth]{./images/english/imf_decomposition.png}};
    \end{tikzpicture}
    \caption{MSCI World closing prices: CEEMDAN versus EEMD when the residue is excluded from the IMF sum. \textbf{Top:} original series and sum of the eight CEEMDAN IMFs (residue excluded). \textbf{Bottom:} $|x(t)-\sum_k \mathrm{IMF}_k(t)|$ without the residue for CEEMDAN ($K=8$) and for EEMD on the same series. The CEEMDAN trace is an almost smooth, monotonically increasing trend (magnitude of the CEEMDAN residue), whereas the EEMD trace is rougher and exhibits fine-scale undulations, indicating residual high-frequency content in the slow component under input-only ensemble averaging.}
    \label{fig:imf_decomposition}
\end{figure}

\FloatBarrier

\subsection{Graph transformation}

Graph transformations were performed via visibility and recurrence algorithms for all IMFs (and the residue) produced by the CEEMDAN decomposition above.

To visually illustrate the resulting structure of these transformations, Figure~\ref{fig:hvg_imf} shows the HVG graph obtained from IMF$_1$, which represents the highest frequency component of the decomposition. The visualization reveals the characteristic topology of HVG graphs: a sparse structure where each node (corresponding to a temporal point of the series) connects only with those nodes that satisfy the horizontal visibility criterion. This structure reflects the high-frequency nature of IMF$_1$, where rapid oscillations generate localized connectivity patterns that preserve the temporal information of the original series. The graphical representation allows us to appreciate how the visibility algorithm captures the geometric relationships between data points, transforming the one-dimensional temporal sequence into a two-dimensional network structure that maintains the structural properties of the modal component.

\begin{figure}[H]
    \centering
    \includegraphics[width=0.8\textwidth]{./images/english/hvg_imf.png}
    \caption{Horizontal visibility graph resulting from the transformation of IMF$_1$ of the MSCI World index. Each node represents a temporal point of the series, and the connections reflect the horizontal visibility criterion.}
    \label{fig:hvg_imf}
\end{figure}

\FloatBarrier

Specifically for the construction of recurrence graphs, optimal embedding parameters were determined independently for each IMF using the methods described in Section~\ref{sec:methodology}. The delay parameter $\tau$ was selected through mutual information analysis, identifying the first minimum of the mutual information function between the original series and its time-delayed version. The embedding dimension $d$ was estimated using the false nearest neighbors (FNN) method, systematically increasing the dimension until the proportion of false neighbors fell below the 10\% threshold. The recurrence threshold $\varepsilon$ was established as the 10th percentile of the distribution of pairwise Euclidean distances between all embedding vectors, ensuring that the graph captures the most significant recurrent structures. The parameter selection results reveal considerable variability among the different IMFs: several high-frequency IMFs select $d=4$, while progressively lower-frequency components tend to reduce the embedding dimension and increase the delay. Table~\ref{tab:parametros_recurrencia} presents the selected parameters for each IMF and the residue.

\begin{table}[h]
\centering
\caption{Embedding parameters and threshold $\varepsilon$ by IMF}
\label{tab:parametros_recurrencia}
\small
\begin{tabular}{lccc}
\toprule
\textbf{Component} & \textbf{$\tau$} & \textbf{$d$} & \textbf{$\varepsilon$} \\
\midrule
IMF$_1$ & 9 & 4 & 0.1847 \\
IMF$_2$ & 2 & 4 & 0.1561 \\
IMF$_3$ & 3 & 4 & 0.1348 \\
IMF$_4$ & 7 & 4 & 0.1887 \\
IMF$_5$ & 15 & 4 & 0.2562 \\
IMF$_6$ & 29 & 3 & 0.1891 \\
IMF$_7$ & 34 & 2 & 0.0798 \\
IMF$_8$ & 32 & 2 & 0.1216 \\
Residue & 23 & 2 & 0.1273 \\
\bottomrule
\end{tabular}
\end{table}

\subsubsection{General structural properties. Connectivity and distance}

The comparative analysis of the fundamental structural properties of graphs generated through the three transformation methods (NVG, HVG, and recurrence) reveals significant differences in the topology of the resulting networks. These differences reflect the inherent characteristics of each transformation algorithm and its interaction with the dynamic properties of the different IMFs.

In terms of network size, visibility-based graphs are built on the full daily grid (3587 nodes for the analyzed closing-price sample). Recurrence graphs are constructed on embedded state vectors, so the number of nodes is slightly smaller (roughly between 3529 and 3581 nodes in this study, depending on $(\tau,d)$). This difference is expected and does not prevent a meaningful comparison of qualitative structural patterns across methods.

The most distinctive characteristic among transformation methods is connection density, measured through the number of links and graph density. HVG graphs are extremely sparse, with densities between approximately $0.00057$ and $0.00111$, corresponding to between 3642 and 7161 undirected edges in the reported construction. Recurrence graphs have moderate densities between approximately $0.00255$ and $0.00287$, with between about 16188 and 18153 undirected edges under the 10th-percentile distance threshold used here. NVG graphs span the widest density range, from about $0.00166$ (high-frequency IMFs) up to about $0.766$ for the trend-dominated residue, with between about 10695 and 4927305 undirected edges depending on the IMF.

The global connectivity of graphs, measured through the number of connected components, reveals another fundamental difference among methods. Both HVG and NVG graphs are connected in this experiment (a single component), which facilitates message passing in many GNN pipelines. Recurrence graphs can fragment: the number of connected components ranges from 1 (notably for the residue in this run) up to about 1590, depending on the IMF. This fragmentation motivates strategies such as pooling over connected components or learned readouts that do not assume a single connected graph.

The average degree follows the same qualitative ordering: HVG graphs have low average degrees (about 2.0--4.0), recurrence graphs have roughly stable moderate average degrees (about 9.1--10.2 under the present thresholding), and NVG graphs vary widely (from single digits up to thousands for the densest modes), reflecting strong scale dependence in natural visibility.

In summary, the three transformations generate graph structures with distinctive topological properties. On one hand, HVG graphs are extremely sparse and always connected, while recurrence graphs present moderate densities but multiple disconnected components. Finally, NVG graphs show highly variable densities that depend strongly on IMF order. These structural differences have important implications for the design of GNN architectures and the selection of appropriate modeling strategies for each type of transformation. Table~\ref{tab:propiedades_estructurales} presents a summary of all structural properties obtained on the graphs after their transformation.

\begin{table}[h]
\centering
\caption{Summary of structural properties of transformed graphs}
\label{tab:propiedades_estructurales}
\small
\begin{tabular}{lccc}
\toprule
\textbf{Property} & \textbf{HVG} & \textbf{Recurrence} & \textbf{NVG} \\
\midrule
Density (range) & $0.00057$ - $0.00111$ & $0.00255$ - $0.00287$ & $0.00166$ - $0.766$ \\
Number of links (range) & $3642$ - $7161$ & $16188$ - $18153$ & $10695$ - $4927305$ \\
Connected components & $1$ (always) & $1$ - $1590$ & $1$ (always) \\
Diameter (range) & $24$ - $3558$ & $1725$ (single-component case) & $10$ - $242$ \\
Average degree (range) & $2.03$ - $3.99$ & $9.1$ - $10.2$ & $6.0$ - $2746$ \\
\bottomrule
\end{tabular}
\end{table}

\subsubsection{Clustering coefficients and local structure}

The clustering coefficient, which quantifies the tendency of nodes to form triangles or local cliques, addresses the question of how cohesive local neighborhoods are, revealing notable contrasts among transformation methods.

A direct comparison of the \emph{average} clustering coefficient ranges in Table~\ref{tab:clustering_coefficients} places NVG highest among the three (across IMF$_1$--IMF$_8$), followed by recurrence graphs, and finally HVG graphs---with the important caveat that NVG becomes extremely dense for the residue (cf.\ Table~\ref{tab:propiedades_estructurales}), which mechanically changes triangle counts, and the global clustering statistics for that NVG were not computed here because of cost.

However, this general hierarchy conceals subtler patterns that depend on IMF order under the CEEMDAN-based decomposition. For NVG graphs on IMF$_1$--IMF$_8$, the average local clustering coefficient ranges from about $0.483$ (IMF$_7$) to $0.754$ (IMF$_2$). The largest values occur at IMF$_2$ and IMF$_1$ ($\approx 0.752$--$0.754$), while IMF$_1$--IMF$_4$ span roughly $0.653$--$0.754$, indicating tight local neighborhoods at high frequency. From IMF$_5$ onward averages drift downward to a minimum near IMF$_7$, before rising again for IMF$_8$ ($\approx 0.607$), so the dependence on IMF index is not monotone. The residue NVG ($\approx 4.9 \times 10^6$ undirected edges) is omitted from the NVG ranges in Table~\ref{tab:clustering_coefficients}, and its regime is dominated by massive local connectivity.

The behavior of HVG graphs contrasts with NVG. IMF$_1$--IMF$_4$ retain comparatively high average clustering (about $0.529$--$0.658$) despite the horizontal-visibility restriction. From IMF$_5$ to IMF$_8$ averages decrease monotonically to about $0.389$. The residue HVG drops to about $0.008$, i.e.\ negligible triangle participation in the mean, consistent with a near-chain visibility pattern on a smooth, slowly varying component---not a complete absence of local clustering at every node, but a strong shift in the distribution. Standard deviations of the node-level clustering coefficients span about $0.065$--$0.289$, largest for the highest-frequency IMFs. Medians range from $0$ (residue) to about $0.667$ (IMF$_1$--IMF$_2$), with several intermediate IMFs at $0.5$, indicating a mixture of nodes with zero local clustering and nodes with moderate clustering.

Recurrence graphs do not follow the same IMF-order trend as HVG. The lowest average clustering appears for IMF$_3$ and IMF$_4$ (about $0.211$--$0.221$), while IMF$_6$, IMF$_8$, and the residue reach the upper part of the range (about $0.409$--$0.459$). This pattern suggests that, under the chosen recurrence construction, mid-frequency embeddings can be the least triangle-rich on average, whereas slower oscillatory modes and the residue align with higher mean local clustering. Standard deviations of the per-node coefficients lie between about $0.254$ and $0.326$. Medians vary widely: some IMFs exhibit a median of $0$ (many embedded states carry zero local clustering), while others reach medians up to about $0.639$.

Table~\ref{tab:clustering_coefficients} synthesizes these results, showing the ranges of average, median, and standard deviation of the \emph{node-level} clustering coefficients for each method (NVG ranges use IMF$_1$--IMF$_8$ only, with the residue NVG omitted).

\begin{table}[h]
\centering
\caption{Clustering coefficients by transformation method and IMF order}
\label{tab:clustering_coefficients}
\small
\begin{tabular}{lccc}
\toprule
\textbf{Method} & \textbf{Average coefficient} & \textbf{Median coefficient} & \textbf{Standard deviation} \\
 & \textbf{(range)} & \textbf{(range)} & \textbf{(range)} \\
\midrule
HVG & $0.0083$ - $0.6584$ & $0.0000$ - $0.6667$ & $0.0652$ - $0.2895$ \\
Recurrence & $0.2109$ - $0.4585$ & $0.0000$ - $0.6389$ & $0.2536$ - $0.3262$ \\
NVG & $0.4832$ - $0.7536$ & $0.4991$ - $0.8333$ & $0.1937$ - $0.3074$ \\
\bottomrule
\end{tabular}
\end{table}

The local structure of transformed graphs, characterized through clustering coefficients, reveals distinctive topological properties that have important implications for GNN modeling. NVG graphs, with high average clustering on oscillatory IMFs (but an omitted residue NVG here), present locally cohesive neighborhoods where message passing can emphasize short-range structure. HVG graphs combine strong clustering at high frequency with a sharp collapse toward near-zero mean clustering on the residue. Recurrence graphs yield intermediate averages with a non-monotone IMF dependence, so modeling choices should reflect embedding dimension, thresholding, and possible fragmentation (as in the connectivity discussion above) alongside these clustering summaries.

\subsubsection{Centrality analysis}

Identifying which nodes play critical roles in the graph structure is fundamental for understanding how information flows and where influence is concentrated. For this purpose, we examine three centrality metrics that capture complementary aspects: betweenness centrality identifies strategic bridges, closeness centrality reveals nodes with rapid access to the rest of the network, and eigenvector centrality measures influence propagated through connections. Unless stated otherwise, summaries refer to the mean of the per-node centralities (NetworkX definitions, normalized betweenness). For recurrence graphs with many edges, betweenness and closeness use the library's sampling-based estimates. For dense NVGs, betweenness uses random $k$-node sampling and closeness uses a random subset of source nodes, both with fixed seeds in the reproducibility script. NVG edge lists for IMF$_1$--IMF$_8$ are loaded from a prior export to avoid rebuilding visibility graphs when possible. The residue NVG is omitted from centrality summaries here because it exceeds the edge budget used in the pipeline.

When analyzing which nodes act as critical bridges, contrasts emerge among transformation methods. \textbf{Betweenness} means stay low on NVGs over IMF$_1$--IMF$_8$ (about $0.0010$--$0.0198$), consistent with dense graphs that spread shortest-path traffic across many nodes. Node-level maxima can still reach about $0.60$ on the highest-frequency NVGs. HVGs show a much wider mean range (about $0.0028$--$0.333$). IMF$_1$ remains near $0.0028$, whereas the residue reaches about $0.333$, reflecting a near-path visibility structure where a handful of dates can act as bottlenecks. Maximum per-node betweenness on HVGs peaks near $0.63$--$0.65$ on several IMFs (including IMF$_1$ and IMF$_8$), far above typical NVG maxima on comparable scales. Recurrence graphs span about $0.0003$--$0.173$ in mean betweenness. The smallest means occur for mid-frequency embeddings (e.g.\ IMF$_3$), while the residue attains the largest mean in this run, consistent with a single large component where a few regions still carry disproportionate inter-cluster traffic under approximate estimation.

\textbf{Closeness} means are highest for NVGs on IMF$_1$--IMF$_8$ (about $0.020$--$0.21$ in this computation, with IMF$_4$ near the upper end and the most oscillatory IMF$_8$ NVG near the lower end), reflecting short average distances in dense visibility graphs. HVG closeness means range from about $0.0009$ on the residue to about $0.091$ on IMF$_1$, while recurrence graphs span about $0.0017$--$0.082$, with IMF$_2$ toward the top and the residue toward the bottom. These patterns align with fragmentation and long effective distances in some recurrence constructions.

\textbf{Eigenvector} centrality means remain modest for HVGs (about $0.0013$--$0.0059$) and recurrence graphs (about $0.0029$--$0.011$), whereas NVG means lie slightly higher (about $0.0051$--$0.011$). Maxima tell a complementary story. HVG eigenvector peaks near $0.62$ on IMF$_8$ despite a low mean, highlighting a few locally influential dates on an otherwise sparse graph. NVG maxima are smaller in relative terms (often below $0.40$ on IMF$_1$--IMF$_3$), while recurrence maxima stay below about $0.15$ except where a dominant component concentrates weight.

\begin{table}[h]
\centering
\caption{Centrality measures by transformation method}
\label{tab:centralidad}
\small
\begin{tabular}{lccc}
\toprule
\textbf{Metric} & \textbf{HVG} & \textbf{Recurrence} & \textbf{NVG} \\
 & \textbf{(average range)} & \textbf{(average range)} & \textbf{(average range)} \\
\midrule
Betweenness & $0.0028$ - $0.3332$ & $0.0003$ - $0.1732$ & $0.0010$ - $0.0198$ \\
Closeness & $0.0009$ - $0.0914$ & $0.0017$ - $0.0815$ & $0.0203$ - $0.2085$ \\
Eigenvector & $0.0013$ - $0.0059$ & $0.0029$ - $0.0107$ & $0.0051$ - $0.0113$ \\
\bottomrule
\end{tabular}
\end{table}

Table~\ref{tab:centralidad} synthesizes the observed ranges of \emph{mean} per-node centralities for each method (NVG column: IMF$_1$--IMF$_8$ only, with the residue NVG omitted). These differences matter for GNN design. NVGs support diffuse, short-range summaries. HVGs can concentrate bridging on a narrow temporal backbone, especially on the residue. Recurrence graphs mix low global closeness with occasional high mean betweenness on the residue, so readouts should tolerate disconnectedness and approximation noise in the chosen estimators.

\subsubsection{Patterns according to IMF order}

Under the CEEMDAN decomposition used here (IMF$_1$--IMF$_8$ plus the residue), IMF index orders components from faster oscillations toward the slow trend. Even so, the graph statistics reviewed above do not track that index in a simple monotone way. Each construction folds scale-dependent dynamics into a different topology, which means ``high-frequency'' and ``low-frequency'' are helpful labels, yet they cannot substitute for component-wise inspection.

Consider first the \textbf{HVG} graphs. From IMF$_1$ toward IMF$_8$, diameters tend to grow and mean clustering drifts downward (see the clustering discussion and Table~\ref{tab:clustering_coefficients}), whereas connectivity remains trivial in the sense of a single component throughout (Table~\ref{tab:propiedades_estructurales}). Against that gradual trend, the residue stands apart: diameter explodes, mean clustering is essentially negligible, and mean betweenness rises sharply compared with the early IMFs (Table~\ref{tab:centralidad}), which is what one expects when horizontal visibility sees a smooth series as an almost one-dimensional chain with a handful of bottleneck dates. For practitioners, the upshot is a clear shift from fairly compact, triangle-rich HVGs at high frequency to a sparse, elongated graph on the residue where bridging can concentrate on very few nodes.

\textbf{Recurrence graphs} paint a contrasting picture. Here, fragmentation is tightly tied to the IMF-specific triple $(\tau,d,\varepsilon)$ in Table~\ref{tab:parametros_recurrencia}. Many modes split into hundreds or more than a thousand components, yet the residue may still appear connected under the same recipe (Table~\ref{tab:propiedades_estructurales}). Nor does clustering follow IMF order in lockstep. Mid-frequency embeddings can be the least triangle-rich on average, while slower oscillations or the residue may sit higher (clustering discussion above). Centrality behaves similarly. Mean betweenness is modest on several mid-frequency IMFs but peaks on the residue in this run (Table~\ref{tab:centralidad}). In practice, it is safer to treat recurrence graphs as a \emph{family} of regimes than as one uniformly ``cohesive'' class. Whenever components multiply, encoders or pooling that work per piece are natural, and any summary that presumes a single connected graph should be checked mode by mode.

Turning finally to \textbf{natural visibility graphs (NVG)}, they exhibit the widest cross-IMF spread in density, degree, diameter, and clustering (Tables~\ref{tab:propiedades_estructurales}--\ref{tab:clustering_coefficients}). IMF$_1$--IMF$_4$ pair comparatively moderate edge counts with strong mean clustering, while deeper IMFs, including IMF$_8$, sit in an intermediate regime (still dense, but far from the residue). At the bottom of the scale, the residue NVG is orders of magnitude heavier than any IMF (Table~\ref{tab:propiedades_estructurales}), to the point that several global summaries were skipped earlier for computational reasons. That limitation is part of the modeling problem itself whenever one insists on a full NVG at the lowest frequency.

In summary, IMF order is best read as a guide to \emph{interpretation}, not as a universal knob. HVG compactness tends to erode toward the residue. Recurrence links embedding choices to fragmentation and to local structure that can rise or fall in ways that are not monotone in IMF index. NVG couples scale to explosive variation in density. For graph-based learning, a workable rule of thumb is that \emph{early IMFs} often pair well with straightforward message passing on connected HVG/NVG views where neighborhoods still carry interpretable structure. \emph{Mid-frequency IMFs} are frequently where recurrence fragmentation and non-monotone clustering force explicit choices about components and thresholds. \emph{The residue} is where path-like HVGs, possibly dominant recurrence betweenness when the graph is connected, and prohibitively dense NVGs all meet, so sparsification, sampling, or multi-scale pooling stop being optional refinements and become part of the design.

\subsubsection{Implications for GNN modeling}

The structural properties identified in transformed graphs have important implications for the application of graph neural networks. NVG graphs, with their high density and very short-distance structure, provide a rich environment for message passing in GNNs, allowing information to propagate efficiently through the network. However, high density can also increase computational complexity and potentially introduce noise in learned representations.

HVG graphs behave differently: their sparser structure yields a leaner representation that reduces computational overhead. Still, the wide range of distances observed—particularly in low-frequency IMFs—can complicate the learning of stable representations. One clear benefit is that all HVG graphs remain connected, enabling global message passing without resorting to specialized techniques for handling isolated components.

Recurrence graphs introduce another layer of complexity. Most IMFs contain numerous disconnected components, forcing GNN designers to adopt component-wise processing or pooling methods that merge information across components. Despite this challenge, recurrence graphs encode core dynamic features of the underlying system, making them useful for prediction tasks.

Taken together, these three methods complement each other. A hybrid approach that integrates multiple graph representations could exploit the advantages of each while compensating for their weaknesses. This direction opens up promising avenues for developing multi-view GNN architectures tailored to financial time series decomposed through CEEMDAN.

\FloatBarrier

\subsection{Dimensionality reduction of the oscillatory IMF panel}
\label{sec:dim_red_imf}

The graph analyses above use the full eight oscillatory IMF channels (plus the residue) on the daily grid, which for some constructions implies a large number of nodes and, especially for NVGs at lower frequency, a very large edge budget. One therefore asks whether those eight channels can be replaced by a smaller set of derived coordinates \emph{at each time index} so that downstream graph construction or repeated experiments could operate on a compact multichannel state. The linear PCA and ICA pipelines used here are specified in Section~\ref{sec:method_dim_red}. This subsection reports only the empirical application to the MSCI CEEMDAN panel.

The same workflow was also run for ranks $k\in\{3,5\}$. Figures and tabulated reconstruction metrics below foreground $k=4$ as a representative mid-rank compromise between compression and fit.

For the MSCI World CEEMDAN panel aligned with the empirical IMF table ($T=3596$ daily observations), the eight oscillatory IMF columns were standardized columnwise and projected to $k=4$ scores using FastICA and, separately, PCA. The CEEMDAN residue was excluded from this linear block so that the slow trend is not folded into the reduced coordinates. Figure~\ref{fig:imf_orig_vs_red} shows the four ICA score series $Z_1$--$Z_4$ over the sample as an illustration of the compact multichannel state. PCA score trajectories can differ at the same rank even when the induced reconstruction agrees with ICA (see below).

\begin{figure}[H]
    \centering
    \includegraphics[width=0.72\textwidth]{./images/english/imf_original_vs_reduced_fastica_k4.png}
    \caption{MSCI World CEEMDAN oscillatory block ($T=3596$): four reduced coordinates $Z_1$--$Z_4$ from ICA after column standardization of IMF$_1$--IMF$_8$ ($k=4$). The CEEMDAN residue is excluded from this linear reduction.}
    \label{fig:imf_orig_vs_red}
\end{figure}

Scores were mapped back through the inverse linear operator in score space and the inverse column standardization so that the fitted oscillatory block $\hat{\mathbf{X}}$ is expressed in the original IMF measurement units. Working in that native domain is deliberate: reconstruction error and information loss after reduction are read by comparing $\hat{\mathbf{X}}$ to $\mathbf{X}$ mode by mode with RMSE and $R^2$ on the CEEMDAN scale rather than only in abstract standardized coordinates. Figure~\ref{fig:imf_orig_vs_recon_panel} plots each $\mathrm{IMF}_j$ alongside $\widehat{\mathrm{IMF}}_j$ from that inverse map. For this sample and rank, ICA and PCA yield \emph{numerically identical} reconstructions $\hat{\mathbf{X}}$, i.e.\ the same rank-four orthogonal projector of the standardized block, while the four score series $Z_j$ need not coincide between methods. In aggregate, $\|\mathbf{X}-\hat{\mathbf{X}}\|_F/\|\mathbf{X}\|_F\approx 0.59$ and the arithmetic mean of the eight columnwise coefficients $R^2_j$ between $\mathrm{IMF}_j$ and $\widehat{\mathrm{IMF}}_j$ is about $0.64$, consistent with the variance captured by the leading four principal directions after column scaling.

\begin{figure}[H]
    \centering
    \includegraphics[width=\textwidth]{./images/english/imf_original_vs_reconstructed_k4.png}
    \caption{MSCI World ($T=3596$). \textbf{Left:} CEEMDAN IMF$_1$--IMF$_8$. \textbf{Right:} rank-$k=4$ linear reconstructions $\widehat{\mathrm{IMF}}_j$ from the standardized oscillatory block (residue excluded).}
    \label{fig:imf_orig_vs_recon_panel}
\end{figure}

Table~\ref{tab:imf_red_recon_k4} reports RMSE and $R^2$ between each native mode and its reconstruction.

\begin{table}[h]
\centering
\caption{Rank-$k=4$ linear reconstruction of each oscillatory IMF on the MSCI sample ($T=3596$): RMSE and $R^2$ between $\mathrm{IMF}_j$ and $\widehat{\mathrm{IMF}}_j$ (ICA and PCA coincide numerically for $\widehat{\mathrm{IMF}}_j$).}
\label{tab:imf_red_recon_k4}
\small
\begin{tabular}{lcc}
\toprule
\textbf{Mode} & \textbf{RMSE} & \textbf{$R^2$} \\
\midrule
IMF$_1$ & 0.380 & 0.300 \\
IMF$_2$ & 0.227 & 0.639 \\
IMF$_3$ & 0.319 & 0.663 \\
IMF$_4$ & 0.488 & 0.747 \\
IMF$_5$ & 0.678 & 0.724 \\
IMF$_6$ & 0.811 & 0.724 \\
IMF$_7$ & 1.789 & 0.642 \\
IMF$_8$ & 4.876 & 0.650 \\
\bottomrule
\end{tabular}
\end{table}

IMF$_1$ shows the weakest $R^2$ (about $0.30$, with RMSE $\approx 0.38$), so most of its variance lies outside the retained four-dimensional subspace. Mid-frequency modes reach $R^2$ up to roughly $0.75$ (IMF$_4$) with RMSE below $0.7$ through IMF$_6$. For IMF$_7$ and IMF$_8$, $R^2$ remains near $0.64$--$0.65$ but RMSE rises sharply ($\approx 1.8$ and $\approx 4.9$), reflecting wider amplitude on the slowest oscillations. Here absolute pointwise errors can be large while the linear fit stays only moderate on a relative scale.

Section~\ref{sec:synthetic_emd} already measures cross-IMF linear coupling through $\bar{\rho}$ and $\pi_{0.05}$ (Table~\ref{tab:synthetic_emd_summary}), partly to justify CEEMDAN as yielding comparatively modest average Pearson redundancy among native modes on the calibrated stress test. Compressing eight IMF channels to four scores $Z_j$ breaks the one-to-one alignment with mode indices, so it is natural to ask whether the $Z_j$ inherit a similar qualitative property: limited mutual linear dependence across parallel channels rather than a few scores that move together like a low-rank copy of the same component. Let $r_{ij}$ denote the sample Pearson correlation between $Z_i$ and $Z_j$ over the $T$ time points, for $1\le i<j\le 4$. For the standardized scores from the PCA and ICA fits above, linear decorrelation of the $Z_j$ is the design target. Empirically, all $|r_{ij}|$ are below $10^{-12}$ (numerical roundoff only). Using the same uncorrected Pearson test protocol as in Table~\ref{tab:synthetic_emd_summary}, none of the six pairs yields $p<0.05$, so $\pi_{0.05}=0$ for this diagnostic. Spearman correlations among the $Z_j$ are not constrained to vanish when Pearson correlations do. Their absolute values average about $0.035$ and reach a maximum near $0.07$, indicating only weak monotonic co-movement while linear decorrelation remains essentially perfect. Higher-order dependence is therefore only partially controlled relative to the ideal ICA objective.

\FloatBarrier

\section{Conclusions and future lines}
\label{conclfutu}

\subsection{Conclusions}
As a result of this study, a methodological framework has been developed that combines EMD with graph transformation techniques to analyze financial time series, specifically applied to the MSCI World index. The results obtained show that this integration of complementary techniques enables a deeper understanding of the complex, nonlinear, and non-stationary dynamics that characterize financial markets.

Statistical validation using the ADF, KPSS, and BDS tests confirmed that the MSCI World index satisfies the necessary conditions to apply EMD effectively, establishing a solid foundation for subsequent analysis. A controlled comparison of EMD, EEMD, and CEEMDAN on synthetic signals (Section~\ref{sec:synthetic_emd}), including a multi-component benchmark with an explicit upward drift and ensemble hyperparameters tuned on that benchmark, illustrated separation limits, mode mixing, and the practical differences between input-only noise injection and stage-wise complete ensemble schemes, motivating the use of CEEMDAN on empirical data. Decomposition using CEEMDAN produced eight IMFs plus a residue, each capturing different temporal scales of market dynamics. Figure~\ref{fig:imf_decomposition} contrasts CEEMDAN with EEMD on the same prices and shows that, when oscillatory IMFs are summed without the residue, CEEMDAN isolates a smooth monotone trend in the implied gap, whereas EEMD leaves noticeable short-term ripple in that quantity, supporting the choice of CEEMDAN for the graph stage. Most energy is concentrated in the residue (slow trend of the price level), while the eight IMFs describe the oscillatory structure around that trend at increasing scales. A complementary linear reduction (Section~\ref{sec:method_dim_red} and empirical results in Section~\ref{sec:dim_red_imf}), positioned after the graph discussion to link compression with large node and edge regimes, illustrated ICA- and PCA-based compression of the eight oscillatory IMF channels to $k\in\{3,4,5\}$ scores per time index, with per-mode reconstruction error and cross-score dependence assessed analogously to the synthetic IMF diagnostics. The empirical graph stage used the native IMF panel without applying that reduction.

The transformation of these IMFs into graph representations through three complementary methods reveals distinctive topological properties that reflect different aspects of the underlying structure of financial data. These properties can be summarized as follows:
\begin{itemize}
    \item HVG graphs generate extremely sparse structures with low density and average degree while maintaining full connectivity. However, they exhibit high variability in distances for low-frequency IMFs, which complicates the learning of consistent representations. 
    \item Recurrence networks show moderate edge densities in the connected regime, but often split into many disconnected components, with counts varying widely across IMFs under the chosen threshold. This fragmentation requires architecture choices (e.g.\ per-component encoders or pooling) that do not assume a single connected graph. 
    \item NVG graphs exhibit the greatest variability in density, ranging from sparse structures for high-frequency IMFs to extremely dense graphs for low-frequency IMFs. This high density, while capturing rich relationships, significantly increases computational complexity.
\end{itemize}


A key finding is that graph statistics relate to IMF scale, yet they need not track IMF index monotonically within a given transformation. HVGs tend toward degraded compactness and very low clustering on the residue. Recurrence graphs couple embedding parameters to fragmentation and IMF-dependent local structure. NVGs exhibit strong scale-dependent density, with an extremely heavy residue graph. This diversity suggests that no single graph representation is uniformly optimal across IMF components.

The results of this study provide crucial insights for the design of GNN architectures applied to financial time series forecasting. The identified topological properties (extreme sparsity in HVGs, disconnected components in recurrence networks, and high density in NVGs) pose specific modeling challenges that must be addressed through specialized architectural strategies. In particular, hybrid approaches are required that can leverage the complementary strengths of multiple graph representations simultaneously.

In conclusion, this work establishes that the integration of EMD with graph analysis constitutes a powerful tool for unraveling the multilevel structure of financial time series. The discovered topological patterns not only deepen our understanding of market dynamics, but also inform the development of more sophisticated GNN-based forecasting systems. The proposed methodology is generalizable to other financial indices and asset classes, opening promising avenues for future research in quantitative financial market analysis.

\subsection{Future lines}
Limitations surfaced here double as a roadmap. One priority is to hard-wire GNN architectures around the topological signatures observed for each IMF graph family. NVG, HVG, and recurrence stress different geometry, so multi-view stacks that ingest several graph views jointly \citep{liu2023multimodal, wang2023hatr} look especially relevant. Adaptive attention could then reweight views as the forecast horizon or volatility regime shifts \citep{su2024attention}.

Recent work already couples temporal encoders with graph convolutions to track both intra-series evolution and cross-asset structure \citep{liu2023multimodal, wang2023hatr, jin2023temporal}. IMF panels add another axis: edges should arguably move with regimes rather than stay static. Recurrence plots here hint at higher-order dependence across modes, so hypergraph constructions \citep{su2024attention, agarwal2022think} deserve prototypes beyond pairwise edges. Concrete architectural bets worth testing next include hyperbolic hypergraph attention for multiscale IMF hierarchies \citep{agarwal2022think}, heterogeneous graphs that type edges by relation \citep{jin2023temporal}, and dual attention over nodes versus relation types to tame mixed signals \citep{wang2023hatr}.

Scalability matters once the universe of assets and refresh cadence grows. Distributed graph engines with partitioning, multilevel parallelism, and communication-aware schedulers already underpin large GNN training \citep{xu2025capturing}. Translating those recipes to CEEMDAN-driven IMF graphs will require partitions that honor temporal adjacency, bounded halo exchange between workers, batched windows for data parallelism, and caches for quasi-stable IMFs that change slowly.

Streams arrive continuously, so graphs built on rolling CEEMDAN outputs need maintenance policies that combine incremental sifting updates instead of full recomputes whenever prices tick, dynamic-graph GNN layers that accept edge insertions and deletions \citep{gravina2023scalable, liu2023decoupled}, and latency-aware serving stacks for live desks. Dense NVGs and fragmented recurrence graphs also stress memory, and sparsifying transforms, neighborhood-importance sampling, and asynchronous training collectives belong on the same roadmap. When densification starts before message passing, widening the pilot in Section~\ref{sec:method_dim_red} beyond linear PCA and FastICA is worth exploring through nonlinear autoencoders, manifold embeddings, sparse factorizations, or kernel lifts that could shrink the channel count before visibility or recurrence rules fire, trading reconstruction bias for thinner adjacency lists under fixed constructors.

Multi-index portfolios compound every scaling issue. Federated workflows could train shared encoders without centralizing order flows. Transfer models fit on liquid benchmarks might warm-start regional books. Aggregation layers could couple macro indices with sector peers instead of treating MSCI World in isolation.

Other research directions include:
\begin{itemize}
    \item Port the methodological framework to additional instruments and check whether the topological motifs generalize,
    \item Track how IMF graph statistics shift across market states and whether those shifts anticipate regime changes,
    \item Build post-hoc or intrinsic interpretability tools that attribute predictions to specific IMFs and graph motifs,
    \item Automate EEMD and CEEMDAN hyperparameter choice (ensemble size, noise scales) per series, extending the grid-search discipline used on the synthetic stress test,
    \item Compare further time-series-to-graph maps and fold in non-price channels (news, social sentiment) where licensing permits,
    \item Push IMF-panel reduction past the PCA and FastICA baseline in Section~\ref{sec:method_dim_red}: benchmark nonlinear and sparse families jointly with the graph constructor when the goal is to cap node and edge growth on each reduced channel.
\end{itemize}

Together, these threads span specialized graph encoders, systems for streaming and distributed training, and richer data interfaces. Crossing them with adaptive decomposition remains a fertile niche for financial forecasting and for risk workflows that must respect latency, privacy, and reconstruction fidelity at once.

\FloatBarrier

\begin{thebibliography}{99}

\bibitem{bollerslev1986generalized}
Bollerslev, T. (1986).
\newblock Generalized autoregressive conditional heteroskedasticity.
\newblock {\em Journal of Econometrics}, 31(3), 307--327.

\bibitem{brock1996test}
Brock, W.~A., Dechert, W.~D., Scheinkman, J.~A., \& LeBaron, B. (1996).
\newblock A test for independence based on the correlation dimension.
\newblock {\em Econometric Reviews}, 15(3), 197--235.

\bibitem{chen2025identifying}
Zeng, Z., \& Chen, Y. (2025).
\newblock Identifying and forecasting recurrently emerging stock trend structures via rising visibility graphs.
\newblock {\em Forecasting}, 7(2), 26.
\newblock \url{https://doi.org/10.3390/forecast7020026}

\bibitem{dickey1979distribution}
Dickey, D.~A., \& Fuller, W.~A. (1979).
\newblock Distribution of the estimators for autoregressive time series with a unit root.
\newblock {\em Journal of the American Statistical Association}, 74(366a), 427--431.

\bibitem{eckmann1987recurrence}
Eckmann, J.~P., Kamphorst, S.~O., \& Ruelle, D. (1987).
\newblock Recurrence plots of dynamical systems.
\newblock {\em EPL (Europhysics Letters)}, 4(9), 973.

\bibitem{engle1982autoregressive}
Engle, R.~F. (1982).
\newblock Autoregressive conditional heteroscedasticity with estimates of the variance of United Kingdom inflation.
\newblock {\em Econometrica}, 50(4), 987--1007.

\bibitem{han2022multiscale}
Han, M., Fan, Q., \& Ling, G. (2022).
\newblock Multiscale online-horizontal-visibility-graph correlation analysis of financial market.
\newblock {\em Physica A: Statistical Mechanics and Its Applications}, 607, 128195.
\newblock \url{https://doi.org/10.1016/j.physa.2022.128195}

\bibitem{hyvarinen2000independent}
Hyvärinen, A., \& Oja, E. (2000).
\newblock Independent component analysis: algorithms and applications.
\newblock {\em Neural Networks}, 13(4--5), 411--430.
\newblock \url{https://doi.org/10.1016/S0893-6080(00)00026-5}

\bibitem{jolliffe2016principal}
Jolliffe, I.~T., \& Cadima, J. (2016).
\newblock Principal component analysis: a review and recent developments.
\newblock {\em Philosophical Transactions of the Royal Society A}, 374(2065), 20150202.
\newblock \url{https://doi.org/10.1098/rsta.2015.0202}

\bibitem{huang1998empirical}
Huang, N.~E., Shen, Z., Long, S.~R., Wu, M.~C., Shih, H.~H., Zheng, Q., \ldots \& Liu, H.~H. (1998).
\newblock The empirical mode decomposition and the Hilbert spectrum for nonlinear and non-stationary time series analysis.
\newblock {\em Proceedings of the Royal Society of London. Series A: Mathematical, Physical and Engineering Sciences}, 454(1971), 903--995.

\bibitem{kwiatkowski1992testing}
Kwiatkowski, D., Phillips, P.~C., Schmidt, P., \& Shin, Y. (1992).
\newblock Testing the null hypothesis of stationarity against the alternative of a unit root: How sure are we that economic time series have a unit root?
\newblock {\em Journal of Econometrics}, 54(1-3), 159--178.

\bibitem{lacasa2008visibility}
Lacasa, L., Luque, B., Ballesteros, F., Luque, J., \& Nuño, J.~C. (2008).
\newblock From time series to complex networks: The visibility graph.
\newblock {\em Proceedings of the National Academy of Sciences}, 105(13), 4972--4975.

\bibitem{ljung1978measure}
Ljung, G.~M., \& Box, G.~E. (1978).
\newblock On a measure of lack of fit in time series models.
\newblock {\em Biometrika}, 65(2), 297--303.

\bibitem{rilling2008one}
Rilling, G., \& Flandrin, P. (2008).
\newblock One or two frequencies? The empirical mode decomposition answers.
\newblock {\em IEEE Transactions on Signal Processing}, 56(1), 85--95.

\bibitem{luque2009horizontal}
Luque, B., Lacasa, L., Ballesteros, F., \& Luque, J. (2009).
\newblock Horizontal visibility graphs: Exact results for random time series.
\newblock {\em Physical Review E}, 80(4), 046103.

\bibitem{song2024early}
Song, S., \& Li, H. (2024).
\newblock Early warning signals for stock market crashes: empirical and analytical insights utilizing nonlinear methods.
\newblock {\em EPJ Data Science}, 13(1), 16.
\newblock \url{https://doi.org/10.1140/epjds/s13688-024-00457-2}

\bibitem{su2025hybrid}
Su, J., Li, H., Wang, R., Guo, W., Hao, Y., Kurths, J., \& Gao, Z. (2025).
\newblock A hybrid dual-branch model with recurrence plots and transposed transformer for stock trend prediction.
\newblock {\em Chaos: An Interdisciplinary Journal of Nonlinear Science}, 35(1), 013125.
\newblock \url{https://doi.org/10.1063/5.0233275}

\bibitem{wu2009ensemble}
Wu, Z., \& Huang, N.~E. (2009).
\newblock Ensemble empirical mode decomposition: A noise-assisted data analysis method.
\newblock {\em Advances in Adaptive Data Analysis}, 1(1), 1--41.

\bibitem{torres2011complete}
Torres, M.~E., Colominas, M.~A., Schlotthauer, G., \& Flandrin, P. (2011).
\newblock A complete ensemble empirical mode decomposition with adaptive noise.
\newblock {\em 2011 IEEE International Conference on Acoustics, Speech and Signal Processing (ICASSP)}, 4144--4147.

\bibitem{liu2023multimodal}
Liu, R., Liu, H., Huang, H., Yang, J., Gao, Y., Zhu, J., Wang, C., \& Xu, K. (2023).
\newblock Multimodal multiscale dynamic graph convolution networks for stock price prediction.
\newblock {\em Pattern Recognition}, 144, 110211.
\newblock \url{https://doi.org/10.1016/j.patcog.2023.110211}

\bibitem{wang2023hatr}
Wang, H., Wang, T., Li, S., Li, J., \& Zhang, Z. (2023).
\newblock HATR-I: Hierarchical Adaptive Temporal Relational Interaction for Stock Trend Prediction.
\newblock {\em IEEE Transactions on Knowledge and Data Engineering}, 35(7), 7492--7506.
\newblock \url{https://doi.org/10.1109/TKDE.2022.3188320}

\bibitem{su2024attention}
Su, H., Wang, X., Qin, Y., Zhao, X., \& Lin, C. (2024).
\newblock Attention based adaptive spatial--temporal hypergraph convolutional networks for stock price trend prediction.
\newblock {\em Expert Systems with Applications}, 237, 121899.
\newblock \url{https://doi.org/10.1016/j.eswa.2023.121899}

\bibitem{jin2023temporal}
Jin, M., Koh, H.~Y., Wen, Q., Zambon, D., Alippi, C., Webb, G.~I., King, I., \& Pan, S. (2023).
\newblock Temporal and Heterogeneous Graph Neural Network for Financial Time Series Prediction.
\newblock In {\em Proceedings of the 31st ACM International Conference on Information and Knowledge Management} (CIKM '22), 3584--3593.
\newblock \url{https://doi.org/10.1145/3511808.3557089}

\bibitem{agarwal2022think}
Agarwal, S., Sawhney, R., Thakkar, M., \& Gupta, D. (2022).
\newblock THINK: Temporal Hypergraph Hyperbolic Network.
\newblock In {\em Proceedings of the 2022 IEEE International Conference on Data Mining} (ICDM), 1--10.
\newblock \url{https://doi.org/10.1109/ICDM54844.2022.00096}

\bibitem{gravina2023scalable}
Gravina, A., Zambon, D., Bacciu, D., \& Alippi, C. (2023).
\newblock Scalable spatiotemporal graph neural networks.
\newblock In {\em Proceedings of the AAAI Conference on Artificial Intelligence}, 37(6), 7688--7696.

\bibitem{liu2023decoupled}
Liu, Y., Li, K., Xie, Y., Guo, Y., Yan, J., \& Yu, P.~S. (2023).
\newblock Decoupled Graph Neural Networks for Large Dynamic Graphs.
\newblock {\em arXiv preprint arXiv:2305.08273}.
\newblock \url{https://doi.org/10.48550/arXiv.2305.08273}

\bibitem{xu2025capturing}
Xu, Q. (2025).
\newblock Capturing Structural Evolution in Financial Markets with Graph Neural Time Series Models.
\newblock {\em Preprints}, 2025071260.
\newblock \url{https://doi.org/10.20944/preprints202507.1260.v1}

\end{thebibliography}

\end{document}

\usepackage{subcaption}
\usepackage{tikz}
\usepackage{placeins}
\usepackage{float}

% Evitar que figuras/tablas acumuladas pasen al final del documento (p.\ ej.\ tras la bibliografía).
\setcounter{topnumber}{9}
\setcounter{bottomnumber}{9}
\setcounter{totalnumber}{20}
\renewcommand{\topfraction}{0.95}
\renewcommand{\bottomfraction}{0.95}
\renewcommand{\textfraction}{0.07}
\renewcommand{\floatpagefraction}{0.85}

\title{Empirical Mode Decomposition and Graph Transformation of the MSCI World Index: A Multiscale Topological Analysis for Graph Neural Network Modeling}

\author{Agustín Manuel de los Riscos Mayorga}

\date{\today}

\AtBeginDocument{\renewcommand{\abstractname}{}}

\begin{document}

\maketitle

\begin{abstract}
This work targets graph representations that can feed graph neural network (GNN) pipelines and therefore applies Empirical Mode Decomposition (EMD) to the MSCI World index, assigning a graph object to each extracted mode. The exposition begins with controlled synthetic benchmarks, including a multi-component stress test with an equity-like upward drift, where classical EMD is compared with Ensemble EMD (EEMD) and Complete Ensemble EMD with Adaptive Noise (CEEMDAN), using ensemble sizes and noise hyperparameters from a small grid search aimed at low average cross-IMF correlation together with tight numerical reconstruction. That comparison clarifies separation limits, mode mixing, and the trade-offs of noise-assisted schemes and therefore grounds the subsequent choice of CEEMDAN for the empirical series, which yields eight intrinsic mode functions (IMFs) spanning high-frequency structure through slower swings. Once those modes are fixed, three deterministic time-series-to-graph maps (natural visibility (NVG), horizontal visibility (HVG), recurrence graphs) attach a graph instance to each IMF and, when relevant, to the residue. Because each map stresses different geometry, the induced topology co-varies with both IMF scale and construction family. Visibility graphs emphasize amplitude-driven connectivity and strong local clustering, whereas recurrence graphs track embedding-space proximity and often fragment, so graph-level statistics need not move monotonically with IMF order.

Retaining every oscillatory IMF as its own channel can inflate node and edge budgets under several graph constructions, and a smaller per-date multivariate summary could ease downstream graph workloads or repeated multichannel experiments. Accordingly, the text documents an optional linear compression of the eight oscillatory IMF columns to $k\in\{3,4,5\}$ scores per date via PCA and FastICA on the MSCI CEEMDAN table, with figures and tables focused on $k=4$. That block evaluates reconstructions in native IMF units, reports numerical agreement of ICA and PCA on $\hat{\mathbf{X}}$ at that rank, and contrasts linear coupling among scores with the synthetic cross-IMF Pearson diagnostics, whereas the headline graph statistics in the core results still use native IMFs plus residue without applying that shrinkage. Taking all together, these results show pronounced heterogeneity across scales and encoders, so scale-aware GNN pipelines should adapt to the view at hand rather than reuse one graph template across modes.
\end{abstract}

\vspace{0.5cm}

{\normalfont\normalsize\noindent\textbf{Keywords:} Empirical Mode Decomposition, EEMD, CEEMDAN, MSCI World, Time Series-Graph Transformation, Visibility Graphs, Recurrence Graphs, Graph Neural Networks, Principal Component Analysis, FastICA, Dimensionality Reduction, Financial Time Series}

\newpage
\tableofcontents
\newpage

\section{Introduction}

Analyzing financial time series poses specific challenges derived from the market, whose dynamics are typically neither linear nor stationary. Traditional time series decomposition methods, including wavelet transforms and Fourier analysis, are based on predefined basis functions that often prove inadequate for capturing the inherent properties of financial datasets. The \textbf{Empirical Mode Decomposition} (EMD), originally proposed by \citep{huang1998empirical}, has proven to be highly suitable for the analysis of financial time series, as it provides an adaptive and data-driven methodology that operates without making assumptions about the underlying data structure. Unlike traditional techniques, EMD adaptively locates local extrema and constructs envelopes to isolate oscillatory components, making it highly effective for decomposing financial series. Each of the resulting series is characterized by different time-dependent amplitude and frequency values.

The effectiveness of EMD in the analysis of financial time series stems from this technique's ability to adapt to the fundamental properties that define stock market data: \textbf{non-stationarity} resulting from prolonged trends, \textbf{nonlinear dynamics} evident through volatility clustering and regime transitions, and \textbf{variable amplitude and frequency patterns} according to different market states. These characteristics make financial time series particularly suitable for EMD analysis, given that the technique is explicitly designed to isolate intrinsic mode functions (IMFs) from signals exhibiting these characteristics. The \textbf{MSCI World} index, which serves as a global market indicator, fulfills the properties enumerated above, positioning it as an appropriate case study for applying EMD to the decomposition of time series in this sector.

Although EMD provides a robust decomposition framework, the resulting IMFs require advanced modeling techniques to represent structural and temporal relationships. To achieve this objective, this line of work proposes \textbf{graph neural networks} (GNNs), given their ability to effectively learn from graph-structured data and identify both localized and global patterns through message-passing mechanisms. However, to utilize GNNs in the analysis of IMFs, it is essential to represent temporal sequences as nodes and edges through data transformations, while maintaining the temporal order of the series at all times. Embedding each IMF as its own graph on the full daily grid multiplies structural objects with the number of modes and can inflate node and edge budgets for visibility-type constructions, which motivates documenting and testing optional linear compression of the oscillatory IMF block after the core graph statistics (Sections~\ref{sec:method_dim_red} and~\ref{sec:dim_red_imf}).

The central thread is CEEMDAN on MSCI prices together with complementary graph encodings of each extracted mode. Concretely, the manuscript covers the following points.
\begin{itemize}
    \item \textbf{validating the applicability conditions of EMD} for the MSCI World index through statistical tests.
    \item \textbf{comparing EMD, EEMD, and CEEMDAN on synthetic signals} of controlled design (Section~\ref{sec:synthetic_emd}), with ensemble hyperparameters for the noise-assisted methods tuned on the multi-component family so that the numerical experiment matches the stress test used in the figure.
    \item \textbf{decomposing the index} into IMFs and a residue via CEEMDAN, with a characterization of the energy, variance, frequency, and amplitude attributes of each oscillatory mode, and \textbf{a direct comparison} on the same MSCI closing prices between CEEMDAN and EEMD in terms of the reconstruction gap when the residue is excluded from the IMF sum (Figure~\ref{fig:imf_decomposition}), to highlight the improvement of stage-wise adaptive noise over input-only ensemble averaging for isolating the slow trend.
    \item \textbf{converting each IMF (and the residue, where applicable)} into graph representations using NVG, HVG, and recurrence constructions.
    \item \textbf{calculating metrics of the obtained graphs} (clustering coefficients, centrality measures, connectivity patterns) that summarize the structural characteristics of each decomposed component.
    \item \textbf{documenting optional PCA and FastICA compression} of the eight oscillatory IMF channels (Section~\ref{sec:method_dim_red}), with empirical reconstruction and score-dependence diagnostics on the MSCI CEEMDAN panel after those graph analyses (Section~\ref{sec:dim_red_imf}), motivated by large node and edge budgets when every mode is embedded separately, foregrounding $k=4$ among ranks $k\in\{3,4,5\}$, mapping fits back to native IMF units for RMSE and $R^2$ by mode, and contrasting Pearson coupling among reduced scores with the synthetic cross-IMF summaries, while the tabulated graph constructions use native IMF$_1$--IMF$_8$ and the residue without applying that reduction.
\end{itemize}


Beyond this opening section, Section~\ref{sec:methodology} collects data description, EMD theory and variants, applicability checks, NVG/HVG/recurrence graph rules, and optional PCA and FastICA compression of the oscillatory IMF panel with the graph-motivated rationale developed in Section~\ref{sec:method_dim_red}. Section~\ref{sec:results} orders the empirical blocks as follows. First come suitability tests for EMD on MSCI (Section~\ref{sec:suitcondemd}). Next is the synthetic EMD versus EEMD versus CEEMDAN comparison (Section~\ref{sec:synthetic_emd}). The third block reports CEEMDAN on the index with IMF diagnostics. The fourth summarizes graph statistics under the three maps. The fifth presents the MSCI compression study (Section~\ref{sec:dim_red_imf}) after the graph discussion, tying rank-$k$ linear summaries to the observed node and edge burden, with per-mode reconstruction in native IMF units, auxiliary ranks $k\in\{3,5\}$ beside the $k=4$ exhibits, and a Pearson check on reduced scores parallel to the synthetic IMF panel. Section~\ref{conclfutu} closes with synthesis and forward-looking items.
\section{Methodology}
\label{sec:methodology}

\subsection{Description of the data}

The analysis is based on the daily closing prices of the MSCI World index (Figure~\ref{fig:msci}). The analyzed dataset spans from January 12, 2012 to April 20, 2026, and comprises 3587 daily observations. This period covers various market regimes, including bull markets, corrections, and periods of high volatility. The price series ranges from a minimum of 37.73\$ in US dollars to a maximum of 195.79\$, representing a cumulative return of 407.52 \% during the analysis period. Daily returns are calculated as percentage variations, with a mean daily return of 0.0511 \% and a standard deviation of 1.0726 \%. The 30-day moving average volatility is 0.9380 \%, with a maximum value of 5.1228 \% during periods of high volatility.

\begin{figure}[H]
    \centering
    \begin{tikzpicture}
        \node[anchor=south west,inner sep=0] (image) at (0,0) {\includegraphics[width=\textwidth]{./images/english/msci_world_data.png}};
        \begin{scope}[x={(image.south east)},y={(image.north west)}]
            % Etiqueta a) en el panel superior - esquina superior izquierda
            \node[anchor=north west, font=\bfseries, fill=white, inner sep=2pt, opacity=0.8] at (-0.01,0.98) {(a)};
            % Etiqueta b) en el panel medio - esquina superior izquierda
            \node[anchor=north west, font=\bfseries, fill=white, inner sep=2pt, opacity=0.8] at (-0.01,0.65) {(b)};
            % Etiqueta c) en el panel inferior - esquina superior izquierda
            \node[anchor=north west, font=\bfseries, fill=white, inner sep=2pt, opacity=0.8] at (-0.01,0.32) {(c)};
        \end{scope}
    \end{tikzpicture}
    \caption{Analyzed data of the MSCI World. \textbf{a)} Time series trends of closing prices for the MSCI World index from January 2012 to April 2026, accompanied by the 50 and 200-day moving averages that allow identifying short and long-term patterns. \textbf{b)} Daily returns calculated as percentage variations, which capture the daily variability of the index.  \textbf{c)} 30-day moving volatility, which measures the degree of variation of returns in a 30-day moving window. This metric allows identifying periods of higher risk and market stability, showing pronounced peaks during market events such as those observed in 2020.}
    \label{fig:msci} % Para referenciarla en el texto con \ref{fig:mi_imagen}
\end{figure}
\subsection{Empirical Mode Decomposition. Theory and variants}

The EMD decomposition, introduced by \citep{huang1998empirical}, is an adaptive data-driven method for decomposing non-stationary and nonlinear time series into a finite set of oscillatory components called intrinsic mode functions (IMF). Unlike traditional decomposition methods that rely on predefined basis functions (for example, Fourier transforms that use sinusoidal bases or wavelet transforms that use predefined mother wavelets), EMD adaptively identifies local extrema and constructs envelopes to extract oscillatory components at different temporal scales.

The fundamental principle of EMD is the sifting process, which iteratively extracts IMFs by identifying local maxima and minima, constructing upper and lower envelopes through cubic spline interpolation, calculating the mean envelope, and subtracting it from the signal. This process continues until the resulting signal satisfies the IMF conditions: the number of extrema and zero crossings must differ by at most one, and the mean value of the upper and lower envelopes must be zero at any point. This process ends when a stopping criterion is met, typically based on a standard deviation threshold between consecutive sifting iterations.

However, standard EMD has several limitations, among which the mode mixing problem stands out. When decomposing with the original EMD method, oscillations of different scales tend to mix into a single IMF and similar oscillations may be split across several IMFs. This problem is magnified in the presence of intermittent signals and noise, which can make the decomposition unstable and non-unique.

The \textbf{Complete Ensemble Empirical Mode Decomposition with Adaptive Noise (CEEMDAN)} \citep{torres2011complete} addresses mode mixing by injecting white noise in an adaptive way at each stage of the decomposition. Noise realizations are drawn in the ensemble, the corresponding IMFs are averaged across trials, and the procedure continues until the full decomposition—including a residual trend—is obtained. Compared with injecting noise only at the first stage, this construction typically yields more stable IMFs with lower residual noise contamination for a given ensemble size.

To obtain good results in the decomposition of the MSCI World, CEEMDAN is applied with the execution parameters defined in Table~\ref{tab:ceemdan_params}. Those values are specific to the empirical series. The exploratory synthetic experiments in Section~\ref{sec:synthetic_emd} use the ensemble sizes and noise hyperparameters reported there (including a grid search on the multi-component stress test), which need not coincide numerically with Table~\ref{tab:ceemdan_params}.

\begin{table}[h]
\centering
\caption{CEEMDAN decomposition parameters}
\label{tab:ceemdan_params}
\begin{tabular}{ll}
\toprule
\textbf{Parameter} & \textbf{Value} \\
\midrule
Maximum number of IMFs & 14 \\
Ensemble trials & 100 \\
Noise amplitude $\varepsilon$ & 0.05 (5\% of signal's std dev) \\
Pseudorandom seed (noise) & 42 \\
\bottomrule
\end{tabular}
\end{table}

\subsection{Prerequisites for effective EMD application.}

For EMD analysis to be effective, the time series must satisfy the following key conditions:

\begin{itemize}
\item \textbf{Non-stationarity}: EMD decomposition offers optimal performance with non-stationary series, as it can adaptively capture trends and components at different temporal scales. Therefore, the stationarity of the studied MSCI World time interval is checked using the Augmented Dickey-Fuller (ADF) test \citep{dickey1979distribution} and the Kwiatkowski-Phillips-Schmidt-Shin (KPSS) test \citep{kwiatkowski1992testing}.

\item \textbf{Non-linearity}: Applying EMD techniques to the series does not require linearity assumptions. Non-linearity is evaluated through multiple approaches: using the Brock-Dechert-Scheinkman (BDS) test \citep{brock1996test} on the residuals of an AR(1) model, through autocorrelation analysis of squared returns to detect ARCH/GARCH effects \citep{engle1982autoregressive,bollerslev1986generalized}, and finally using the Ljung-Box test \citep{ljung1978measure} on the squared returns of the index.

\item \textbf{Variable amplitude and frequency}: EMD excels at capturing modes with time-varying amplitude and frequency. Therefore, amplitude variability is analyzed through continuous volatility measures (252-day windows corresponding to one trading year) and frequency variability through zero-crossing analysis.
\end{itemize}

The results of all tests listed in this section are described in Section~\ref{sec:suitcondemd}.

\subsection{Time series to graph transformations}

The transformation of time series into graph representations allows us to analyze the structural properties of temporal data using graph theory techniques and network analysis. Furthermore, converting to graphs enables the application of graph neural networks and offers an alternative perspective for understanding the underlying dynamics of decomposed time series components. In this study, we employ three complementary graph transformation methods to convert the extracted IMFs into graph structures: the natural visibility graph (NVG), the horizontal visibility graph (HVG), and recurrence-based networks.

\subsubsection{Visibility algorithms}

Visibility graph algorithms, introduced by \citep{lacasa2008visibility}, transform a time series into a graph by establishing connections between data points based on geometric visibility criteria. The fundamental principle is that two points in the time series are connected if no intermediate point obstructs the line of sight between them.

The variant called the natural visibility graph (NVG) algorithm constructs a graph in which each data point $(t_i, y_i)$ of the time series becomes a node, and an edge exists between nodes $i$ and $j$ if, for all intermediate points $k$ such that $i < k < j$, the following condition is satisfied:
\begin{equation}
y_k < y_i + (y_j - y_i) \frac{t_k - t_i}{t_j - t_i}
\end{equation}

This geometric criterion ensures that the line segment connecting points $i$ and $j$ lies above all intermediate data points, establishing a natural visibility relationship. NVG preserves important properties of the original time series, such as periodicity, fractality, and hierarchical structure, making it suitable for analyzing complex temporal patterns in financial data.

In a second variant, the horizontal visibility graph (HVG) algorithm \citep{luque2009horizontal} simplifies NVG by imposing a stricter horizontal visibility constraint. Two nodes $i$ and $j$ are connected if, for all intermediate points $k$ with $i < k < j$, the following condition is satisfied:
\begin{equation}
y_k < \min(y_i, y_j)
\end{equation}

These two transformations have demonstrated utility in various recent applications in the financial domain. An example is found in Han et al. (2022), where they developed MOHVGCA (Multiscale Online-Horizontal-Visibility-Graph Correlation Analysis), a method that employs HVG to quantify scale-dependent correlations between financial time series, revealing dynamic associations between WTI oil futures and various currencies during oil crises \citep{han2022multiscale}. More recently, Zeng and Chen (2025) proposed the use of rising visibility graphs to extract topological motifs from trend structures in stock indices, demonstrating that these recurrent topological structures can be identified and used to forecast emerging patterns in financial time series \citep{chen2025identifying}. These studies demonstrate the ability of visibility algorithms to reveal hidden structural characteristics in time series from various domains, validating their application in the analysis of complex financial data as addressed in this work.

\subsubsection{Recurrence-based algorithms}

As an alternative to visibility algorithms, the recurrence matrix can be employed as an adjacency matrix to transform time series into graph representations by identifying recurrence patterns within the dynamical system \citep{eckmann1987recurrence}. This method captures the recurrence structure of the series in a higher-dimensional embedding space. This effectively reveals both the geometry of the underlying attractor and the essential dynamical properties of the system.

The construction of recurrence graphs begins with the determination of optimal embedding parameters using geometric and information-theoretic criteria. The delay parameter $\tau$ is selected by analyzing mutual information, specifically identifying the first minimum of the mutual information function between the original series and its time-delayed version. This ensures that the delayed coordinates provide maximum information about the system state while minimizing redundancy. Subsequently, the optimal embedding dimension $d$ is estimated using the false nearest neighbors (FNN) method. This technique identifies the minimum dimension required to unfold the attractor by detecting false neighbors that appear close in low dimensions but are separated in higher dimensions. The FNN method systematically increases the embedding dimension until the proportion of false neighbors falls below a threshold, typically 10\%.

Once $\tau$ (time delay) and $d$ (embedding dimension) are determined, the embedding space is constructed using delay coordinate embedding. Each point $X_i$ in this space is defined as a $d$-dimensional vector:
\begin{equation}
X_i = [Y_i, Y_{i+\tau}, Y_{i+2\tau}, \ldots, Y_{i+(d-1)\tau}]
\end{equation}
where $X_i$ represents the $i$-th point in the embedding space and $Y_i$ denotes the value of the original time series at time $i$.

The resulting recurrence graph preserves the topological and geometric properties of the underlying attractor, enabling the analysis of dynamical invariants such as correlation dimension, Lyapunov exponents, and entropy measures. Node features in the recurrence graph are assigned using the first component of the embedding vectors, maintaining a connection with the original time series values.

From the generated embedding space, the recurrence graph is constructed by computing pairwise distances between all embedding vectors. Two nodes $i$ and $j$ in the graph are connected if the distance between their corresponding embedding vectors $X_i$ and $X_j$ is below a threshold $\varepsilon$. This is formally expressed through the $(i, j)$ element of the recurrence matrix $R$:
\begin{equation}
R_{ij} = \begin{cases}
1 & \text{if } \left\|X_i - X_j\right\| \le \varepsilon \\
0 & \text{otherwise}
\end{cases}
\end{equation}
where $R_{ij}$ is the $(i, j)$-th element of the recurrence matrix. The threshold $\varepsilon$ is typically chosen as a percentile of the pairwise distance distribution, ensuring that the graph captures the most significant recurrent structures while maintaining computational tractability.

Numerous studies have demonstrated that the transformation of time series to recurrence graphs has proven particularly useful in financial market analysis, where capturing nonlinear dynamics and complex patterns is crucial for prediction and anomaly detection. In this context, a hybrid dual-branch architecture combines recurrence plots with transposed transformers for stock trend prediction \citep{su2025hybrid}. In that setup, closing price time series are transformed into recurrence plots using sliding windows of historical segments. The first branch generates RP images from these univariate sequences and extracts recurrence quantification analysis (RQA) measures, which capture nonlinear characteristics and temporal relationships between points. These RQA measures include metrics such as recurrence rate (RR), determinism (DET), and laminarity (LAM), which quantify the predictability and structure of the dynamical system. The second branch processes multivariate time series features using a transposed transformer, and both branches are merged to classify the next-day trend, either up or down. Evaluation on seven selected stocks indicates that combining recurrence-based features with sequence modeling through transformers can yield higher accuracy and F1-scores than traditional machine learning and deep learning baselines. This line of work illustrates how recurrence plots can capture nonlinear dynamics and relationships between temporal points that are difficult to detect using conventional time series analysis methods.

Separately, multiplex recurrence networks (MRNs) have been used to search for early warning signals of critical transitions and crashes in stock markets \citep{song2024early}. In that framework, multidimensional return vectors are transformed into multiplex recurrence networks, where each dimension, corresponding to different assets or indicators, forms a network layer. Recurrence relationships are established within each layer and between layers, capturing both the individual dynamics of each asset and the interdependencies between them. Topological summaries of MRNs, specifically average mutual information and average edge overlap, track changes in the recurrence structure over time. Those summaries are assessed on simulated systems and on empirical data from Chinese and US markets, where indicators derived from MRNs appear to anticipate imminent state transitions and to shed light on major historical market episodes.

Following the review of transformation types, it is concluded that applying visibility and recurrence algorithms to IMF (Intrinsic Mode Function) components is of considerable interest. This methodological strategy enables a more comprehensive analysis of the resulting graph, as visibility complements information regarding the temporal structure of the sequence, while recurrence reveals the relationships of the dynamical system inherent to the time series. As discussed in Section~\ref{conclfutu}, applying other types of transformations will be considered in future research lines.

\subsection{Linear dimensionality reduction of a multichannel IMF panel}
\label{sec:method_dim_red}

An EMD-type decomposition turns one observed series into several aligned oscillatory channels indexed by mode order. That expansion is valuable for scale-resolved diagnostics, but in the graph-oriented workflow emphasized in this article it also multiplies the objects passed to later stages: each IMF can be embedded as its own graph, and aggregate cost then scales with the number of modes as well as with the node and edge burden of the chosen visibility or recurrence construction (Section~\ref{sec:results}). Linear reduction of the IMF panel is treated as an \emph{optional} companion: usefulness is checked empirically rather than postulated. The target is to see whether a handful of derived coordinates per date can carry enough joint IMF variation when multivariate width, repeated experiment cost, or encoder size must be capped, at the cost of losing one-to-one alignment with individual IMFs.

Let $\mathbf{X}\in\mathbb{R}^{T\times p}$ denote a panel of $p$ aligned oscillatory IMFs (columns) over $T$ time indices (rows), obtained after an EMD-type decomposition technique. The aim is to obtain, at each time row and without temporal subsampling, a smaller set of $k<p$ derived coordinates by applying a linear map from $\mathbb{R}^{p}$ to $\mathbb{R}^{k}$ rowwise after a fixed affine preprocessing of the columns. Because this construction linearly recombines IMF amplitudes, it sacrifices the scale-specific interpretation of individual IMFs under EMD. When graph construction or repeated experiments constrain multichannel input size, however, it offers a compact multivariate summary of the oscillatory block. To keep the slow trend from contaminating the reduced coordinates, any EMD residue is excluded from the linear block whenever it is present. To place IMFs on a comparable scale before the rowwise map, the oscillatory columns are standardized (centering and unit-variance scaling) prior to projection, following the affine convention used in the PCA and ICA subsections below.

\subsubsection{Principal component analysis (PCA)}
PCA seeks orthogonal directions of maximal variance after centering the columns (here, after the standardization step above) \citep{jolliffe2016principal}. The resulting score matrix $\mathbf{Z}\in\mathbb{R}^{T\times k}$ has columns that are pairwise uncorrelated by construction and orders directions by explained variance in the standardized feature space. Retaining only $k<p$ components is the best rank-$k$ approximation to the standardized data in the Frobenius norm (equivalently, in the least-squares sense): an orthogonal projection onto a $k$-dimensional subspace. Denoting the reconstructed oscillatory block by $\hat{\mathbf{X}}$, the approximation error $\mathbf{X}-\hat{\mathbf{X}}$ is generally nonzero whenever $k<p$.

\subsubsection{Independent component analysis (ICA)}
ICA postulates a linear mixing model in which each row of the observed (whitened) matrix is a linear combination of latent sources that are mutually \emph{statistically independent} rather than merely uncorrelated \citep{hyvarinen2000independent}. Estimation typically proceeds after whitening (decorrelation and unit variance) and uses measures of non-Gaussianity or contrast functions related to mutual information. FastICA is a standard fixed-point algorithm in this family \citep{hyvarinen2000independent}. When the mixing is invertible and the sources satisfy the usual identifiability conditions, ICA links to blind source separation. In the present setting it serves as a complementary linear reduction whose scores are driven toward independence. For a given $k$, the induced reconstruction of the standardized IMF block need not coincide with that of PCA, although both are linear projections of the same column-standardized panel.

\section{Results}
\label{sec:results}

\subsection{Suitable conditions for EMD application.}
\label{sec:suitcondemd}

A battery of tests is performed to confirm an adequate study via EMD decomposition of the indexed series. Our comprehensive evaluation of the MSCI World index confirms that the time series meets the three fundamental conditions for effective EMD analysis. The results of the three evaluated conditions are presented below.

\subsubsection{Stationarity}
Stationarity tests provide strong evidence of the non-stationarity of the price series, which is ideal for EMD. The Augmented Dickey-Fuller (ADF) test on the closing price series yields a test statistic of 1.6628 with a p-value of 0.9980, which does not allow rejecting the null hypothesis of a unit root at all conventional significance levels. The critical values at the 1\%, 5\% and 10\% levels are -3.4322, -2.8624 and -2.5672, respectively, all substantially more negative than the test statistic. The Kwiatkowski-Phillips-Schmidt-Shin (KPSS) test, which has stationarity as the null hypothesis, produces a test statistic of 8.3015 with a p-value of 0.0100, which rejects the null hypothesis of stationarity at the 5\% significance level. Conversely, daily returns are confirmed as stationary in both tests: ADF statistic of -19.7278 (p-value $<$ 0.0001) and KPSS statistic of 0.0447 (p-value = 0.1000). These complementary results confirm that the price series exhibits non-stationary behavior driven by long-term trends, making it highly suitable for EMD decomposition.

\subsubsection{Linearity}
Multiple tests provide consistent evidence of strong nonlinear dynamics in the MSCI World index. The Brock-Dechert-Scheinkman (BDS) test, applied to the residuals of an AR(1) model, detects significant nonlinear dependencies across multiple embedding dimensions. In dimensions 3, 4 and 5, the BDS test statistics are 20.99, 24.95 and 27.83, respectively, all with p-values $<$ 0.0001, which strongly rejects the null hypothesis of independence and indicates the presence of a nonlinear structure.

The distributional analysis further supports the nonlinear characterization. The price series shows significant deviations from normality, with a skewness of 0.8204 and a kurtosis of -0.2243. The Jarque-Bera test statistic yields a p-value $<$ 0.0001, which strongly rejects the null hypothesis of normality. These findings collectively demonstrate that the MSCI World index exhibits rich nonlinear dynamics, making EMD an adequate decomposition method.

\subsubsection{Volatility in amplitude and frequency}

The autocorrelation analysis of squared returns reveals substantial volatility clustering, a distinctive feature of nonlinear ARCH/GARCH effects. We observe 20 significant autocorrelation lags out of the 20 analyzed, with a maximum autocorrelation of 0.9978, indicating strong volatility persistence. The Ljung-Box test on squared returns produces a large test statistic with a p-value $<$ 0.0001, which provides overwhelming evidence of ARCH/GARCH effects and confirms nonlinear dependence in the variance process.

\subsection{Exploratory comparison of EMD-type methods on synthetic signals}
\label{sec:synthetic_emd}

Before applying decomposition techniques to the MSCI World index, we compared three decomposition families---\textbf{EMD}, \textbf{EEMD} \citep{wu2009ensemble}, and \textbf{CEEMDAN} \citep{torres2011complete}---on \emph{synthetic} time series with known structure. The goal is to stress stylized regimes linked to mode mixing and scale separation: closely spaced harmonics (where amplitude beating stresses frequency-resolution limits \citep{rilling2008one}), linearly swept-frequency chirps (non-stationary instantaneous frequency), intermittent high-frequency bursts on a low-frequency carrier, a compact superposition with additive Gaussian noise, and---as an aggregate stress test---a \emph{multi-component} waveform that sums mode-mixing, ascending and descending chirps, close tones, two sinusoids at $0.5\,\mathrm{Hz}$ and $3\,\mathrm{Hz}$, stronger Gaussian noise, and a slow monotonic uptrend (linear ramp) to mimic the cumulative drift of an equity index. This exercise is illustrative rather than a formal identifiability study. It informs how aggressively each method allocates energy across IMFs when the ground-truth components are not orthogonal harmonics. Figure~\ref{fig:synthetic_signals} shows the five synthetic families (duration $5\,\mathrm{s}$, sampling rate $500\,\mathrm{Hz}$), generated with the same parameter settings as the numerical experiments reported below. The exposition first compares native IMF panels (Table~\ref{tab:synthetic_emd_summary}), then applies FastICA to the same stored decompositions (Table~\ref{tab:synthetic_ica_summary}), and finally states how both blocks guide the CEEMDAN+ICA choice on the empirical index.

\begin{figure}[H]
    \centering
    \includegraphics[width=\textwidth]{./images/english/synthetic_signals_emdsynth.png}
    \caption{Synthetic signals used for the exploratory comparison of EMD-type methods. \textbf{(a)} Sum of two sinusoids with nearby frequencies ($10\,\mathrm{Hz}$ and $11.3\,\mathrm{Hz}$), producing amplitude beating. \textbf{(b)} Linear chirp with initial frequency $1.5\,\mathrm{Hz}$ and sweep rate $1.2\,\mathrm{Hz\,s^{-1}}$. \textbf{(c)} Low-frequency carrier ($0.8\,\mathrm{Hz}$) plus a higher-frequency burst ($12\,\mathrm{Hz}$) active only on $t\in[1.2, 3.2]\,\mathrm{s}$. \textbf{(d)} Superposition of the close-tone signal, $0.35$ times the chirp from panel (b), and additive Gaussian noise scaled to $3\%$ of the standard deviation of the close-tone component. \textbf{(e)} Multi-component benchmark that aggregates the elementary generators: mode mixing ($0.3$--$7\,\mathrm{Hz}$ burst on $t\in[1.25, 3.75]\,\mathrm{s}$), ascending chirp ($f_0=0.8\,\mathrm{Hz}$, $k=0.2\,\mathrm{Hz\,s^{-1}}$), descending chirp ($f_0=4\,\mathrm{Hz}$, $k=-0.15\,\mathrm{Hz\,s^{-1}}$), close tones ($1.5$ and $1.7\,\mathrm{Hz}$), sinusoids at $0.5\,\mathrm{Hz}$ and $3\,\mathrm{Hz}$, additive Gaussian noise with standard deviation $0.2$, and a slow linear uptrend from $0$ to $8$ amplitude units over the $5\,\mathrm{s}$ window.}
    \label{fig:synthetic_signals}
\end{figure}

For each family, the same signal was decomposed with EMD, with EEMD (noise-assisted ensemble averaging at the input), and with CEEMDAN (noise-assisted, stage-wise construction aimed at a \emph{complete} ensemble decomposition \citep{torres2011complete}). Unless stated otherwise, we used a shared cap of $K=12$ intrinsic modes (plus the residue as the final component when applicable). Ensemble sizes were set to $M_{\mathrm{EEMD}}=100$ and $M_{\mathrm{CEEMDAN}}=180$, with CEEMDAN noise scale $\varepsilon=0.03$ and EEMD hyperparameters $(\mathrm{noise\_width},\mathrm{SD\_thresh})=(0.095,0.18)$, values selected by a small grid search on the multi-component benchmark (panel (e)) to reduce average cross-IMF correlation while maintaining numerical reconstruction. For EEMD, the residual supplied with the ensemble mean was included so that the IMFs plus the residual reproduce the input up to numerical roundoff in every case. Empirically, relative reconstruction errors $\|\hat{x}-x\|/\|x\|$ stayed below $10^{-15}$ across all runs.

Table~\ref{tab:synthetic_emd_summary} summarizes linear dependence among \emph{native} intrinsic modes. Let $r_{ij}$ denote the sample Pearson correlation between IMF$_i$ and IMF$_j$ over the $T=2500$ samples, for all pairs among the first $\min\{8,N_{\mathrm{IMF}}\}$ intrinsic modes (excluding the residue). For each pair we test the two-sided hypothesis $H_0:\rho_{ij}=0$ using the standard Pearson correlation test (equivalent to a $t$-statistic for $r_{ij}$ with $T-2$ degrees of freedom under classical Gaussian assumptions). We report $\bar{\rho}=\mathrm{mean}_{i<j}|r_{ij}|$ and the proportion $\pi_{0.05}$ of pairs with $p<0.05$ \emph{without} multiple-comparison adjustment across pairs. With large $T$, even modest $|r_{ij}|$ can yield very small $p$-values (high power), so $\pi_{0.05}$ is interpreted jointly with $\bar{\rho}$ rather than as standalone evidence of ``significance'' in the econometric sense.

\begin{table}[ht!]
\centering
\scriptsize
\setlength{\tabcolsep}{3pt}
\caption{Synthetic experiments: $N$ is the number of returned components (including the residue). $\bar{\rho}$ is the mean absolute Pearson correlation between IMF pairs, and $\pi_{0.05}$ is the fraction of pairs with $p<0.05$ (uncorrected for multiple tests). The number of IMF pairs entering these summaries ranges from $3$ to $28$ depending on $N$ and the eight-mode cap. Same settings as Figure~\ref{fig:synthetic_signals}. The block below the double rule isolates the multi-component stress-test (panel (e)). \textbf{Boldface} in that row marks the method with the smallest $\bar{\rho}$ for that benchmark (weakest average cross-IMF linear coupling).}
\label{tab:synthetic_emd_summary}
\begin{tabular}{l*{9}{c}}
\toprule
 & \multicolumn{3}{c}{\textbf{EMD}} & \multicolumn{3}{c}{\textbf{EEMD}} & \multicolumn{3}{c}{\textbf{CEEMDAN}} \\
\cmidrule(lr){2-4} \cmidrule(lr){5-7} \cmidrule(lr){8-10}
\textbf{Family} & $N$ & $\bar{\rho}$ & $\pi_{0.05}$ & $N$ & $\bar{\rho}$ & $\pi_{0.05}$ & $N$ & $\bar{\rho}$ & $\pi_{0.05}$ \\
\midrule
Close tones & 6 & 0.099 & 0.40 & 12 & 0.072 & 0.21 & 8 & 0.090 & 0.57 \\
Linear chirp & 4 & 0.080 & 0.33 & 12 & 0.089 & 0.29 & 7 & 0.071 & 0.47 \\
Burst on carrier & 4 & 0.071 & 0.67 & 11 & 0.103 & 0.29 & 8 & 0.108 & 0.29 \\
Close tones + chirp + noise & 8 & 0.064 & 0.43 & 12 & 0.092 & 0.29 & 8 & 0.087 & 0.19 \\
\addlinespace[3pt]
\midrule
\midrule
\multicolumn{10}{l}{\scriptsize\emph{Stress-test: highly aggregated multi-component signal (Figure~\ref{fig:synthetic_signals}(e))}} \\[2pt]
\textbf{Multi-comp.\ (panel (e))} & 8 & 0.040 & 0.19 & 11 & 0.071 & 0.32 & \textbf{8} & \textbf{0.038} & \textbf{0.24} \\
\bottomrule
\end{tabular}
\end{table}

Because native IMF columns can retain substantial linear dependence even after a noise-assisted decomposition (Table~\ref{tab:synthetic_emd_summary}), we next apply \textbf{FastICA} \citep{hyvarinen2000independent} to the same stored IMF tables. The oscillatory block uses the leading $p=\min\{8,N_{\mathrm{IMF}}\}$ intrinsic modes in IMF index order (matching the cross-IMF window of Table~\ref{tab:synthetic_emd_summary}); the residue is excluded from the fit. FastICA is run with rank $k=\min\{4,p\}$, paralleling the $k=4$ reduction used later on the MSCI CEEMDAN panel (when $p<4$, $k=p$). Column-standardization precedes ICA, as in Section~\ref{sec:method_dim_red}. Table~\ref{tab:synthetic_ica_summary} reports the same pairwise diagnostics on the $k$ estimated sources; the column $N$ lists the post-reduction channel count $k$ plus the unmodified residue when present. The numerical entries use FastICA with pseudorandom seed $42$, sample length $T=2500$, and output dimension $k=\min\{4,p\}$ as above.

\begin{table}[ht!]
\centering
\scriptsize
\setlength{\tabcolsep}{3pt}
\caption{Post-\textbf{FastICA} \emph{with rank reduction} on the synthetic IMF panels from Table~\ref{tab:synthetic_emd_summary}. The estimator reads $p=\min\{8,N_{\mathrm{IMF}}\}$ IMF columns and outputs $k=\min\{4,p\}$ sources (same $k=4$ target as the MSCI ICA block whenever $p\ge 4$). \textbf{$N$} is the channel count \emph{after} reduction: $N=k+1$ when the residue is appended unchanged, otherwise $N=k$. $\bar{\rho}$ and $\pi_{0.05}$ are computed from Pearson correlations among the $k$ ICA sources only (uncorrected pairs, as in Table~\ref{tab:synthetic_emd_summary}). Table entries follow the protocol above (seed $42$, $T=2500$, $k=\min\{4,p\}$); $\bar{\rho}$ is shown in scientific notation with mantissas rounded to two decimals; $\pi_{0.05}=0$ identically.}
\label{tab:synthetic_ica_summary}
\begin{tabular}{l*{9}{c}}
\toprule
 & \multicolumn{3}{c}{\textbf{EMD}} & \multicolumn{3}{c}{\textbf{EEMD}} & \multicolumn{3}{c}{\textbf{CEEMDAN}} \\
\cmidrule(lr){2-4} \cmidrule(lr){5-7} \cmidrule(lr){8-10}
\textbf{Family} & $N$ & $\bar{\rho}$ & $\pi_{0.05}$ & $N$ & $\bar{\rho}$ & $\pi_{0.05}$ & $N$ & $\bar{\rho}$ & $\pi_{0.05}$ \\
\midrule
Close tones & 5 & $4.79\times 10^{-16}$ & $0$ & 5 & $8.39\times 10^{-16}$ & $0$ & 5 & $5.61\times 10^{-16}$ & $0$ \\
Linear chirp & 4 & $3.67\times 10^{-16}$ & $0$ & 5 & $1.14\times 10^{-14}$ & $0$ & 5 & $1.06\times 10^{-15}$ & $0$ \\
Burst on carrier & 4 & $3.33\times 10^{-16}$ & $0$ & 5 & $1.07\times 10^{-15}$ & $0$ & 5 & $5.64\times 10^{-16}$ & $0$ \\
Close tones + chirp + noise & 5 & $2.97\times 10^{-15}$ & $0$ & 5 & $9.11\times 10^{-14}$ & $0$ & 5 & $6.06\times 10^{-16}$ & $0$ \\
\addlinespace[3pt]
\midrule
\midrule
\multicolumn{10}{l}{\scriptsize\emph{Stress-test: highly aggregated multi-component signal (Figure~\ref{fig:synthetic_signals}(e))}} \\[2pt]
\textbf{Multi-comp.\ (panel (e))} & 5 & $3.58\times 10^{-15}$ & $0$ & 5 & $4.76\times 10^{-16}$ & $0$ & 5 & $6.34\times 10^{-16}$ & $0$ \\
\bottomrule
\end{tabular}
\end{table}

The native-mode statistics are consistent with the intended roles of the three algorithms. Classical EMD returned the smallest $N$ on the linear chirp and burst-on-carrier families ($N=4$ in both cases), but that parsimony aligns with the known mode-mixing risk under intermittency. On the burst family, the variance share of IMF$_1$ was about $0.37$ for EMD versus order $10^{-3}$ for EEMD and order $10^{-4}$ for CEEMDAN, whereas $\bar{\rho}$ ranked highest for CEEMDAN ($0.108$), then EEMD ($0.103$), then EMD ($0.071$). EEMD often hit the $N=12$ cap on simpler families and stopped at $N=11$ on burst and multi-component rows, while CEEMDAN returned $N\in\{7,8\}$ across the five rows (Table~\ref{tab:synthetic_emd_summary}). On the multi-component benchmark (Figure~\ref{fig:synthetic_signals}(e)), with hyperparameters tuned on that panel, CEEMDAN attained the smallest cross-IMF $\bar{\rho}=0.038$ with $N=8$, slightly below EMD ($0.040$), whereas EEMD remained the most redundant at $\bar{\rho}=0.071$ with $N=11$. Residue variance shares on that stress test were about $0.76$ (CEEMDAN), $0.73$ (EMD), and $0.15$ (EEMD), illustrating how the explicit drift is partitioned across IMFs and the residue.

FastICA on the synthetic panels at rank $k=\min\{4,p\}$ on the same $p=\min\{8,N_{\mathrm{IMF}}\}$ IMF window (Table~\ref{tab:synthetic_ica_summary}) drives mean absolute Pearson correlation between the $k$ estimated sources to numerical roundoff and sets $\pi_{0.05}=0$ for every family and every decomposition backend, so the linear coupling visible in Table~\ref{tab:synthetic_emd_summary} is removed in the reduced subspace whether the extractor was EMD, EEMD, or CEEMDAN. For the \emph{empirical} MSCI World pipeline, the controlled evidence therefore motivates \textbf{CEEMDAN} as the decomposition chosen on the drift-rich stress test, together with \textbf{FastICA on the oscillatory IMF block} when downstream stages require both a compact multivariate summary (the table's $N=k+1$ counts $k$ scores plus the residue, in line with rank-$k$ scores when $k<p$ on MSCI) and channels whose \emph{sample} Pearson correlations between rotated components are effectively zero---not merely a parsimonious panel, but a post hoc cancellation of linear redundancy among decomposed modes that CEEMDAN alone does not fully eliminate. Full statistical independence is stronger than zero linear correlation; the ICA table documents the latter at machine precision. IMF indices should still not be read as uniquely identifying independent economic shocks when the data-generating process is unknown.

\FloatBarrier

\subsection{CEEMDAN decomposition of the MSCI World index}
\label{sec:ceemdan_msci}

The decomposition of the MSCI World index was performed using CEEMDAN with the parameters specified in Table~\ref{tab:ceemdan_params}. This choice targets stable mode separation under noise-assisted ensemble averaging, with adaptive noise injection across decomposition stages \citep{torres2011complete}.

The application of CEEMDAN to the closing price series of the MSCI World index, which comprises 3587 daily observations from January 12, 2012 to April 20, 2026, produced a decomposition into \textbf{eight intrinsic mode functions (IMFs)} and a \textbf{residue}. The residue captures the slow-varying level of the index. It is not required to be strictly monotonic in this implementation, but it carries the bulk of the low-frequency energy removed from the oscillatory IMFs.

\subsubsection{Characterization of the IMFs}

Each extracted IMF represents oscillations at different temporal scales, from high-frequency components (IMF$_1$) to lower-frequency components (IMF$_8$). Table~\ref{tab:imf_metrics} summarizes energy, variance, a simple oscillation-count proxy based on local maxima (reported as ``cycles''), and the mean absolute value as an amplitude summary.

\begin{table}[h]
\centering
\caption{Main metrics of the IMFs extracted via CEEMDAN}
\label{tab:imf_metrics}
\small
\begin{tabular}{lcccc}
\toprule
\textbf{IMF} & \textbf{Energy} & \textbf{Variance} & \textbf{Frequency} & \textbf{Amplitude} \\
 & & & \textbf{(cycles)} & \textbf{Average} \\
\midrule
IMF$_1$ & $7.34 \times 10^2$ & 0.2045 & 1287.0 & 0.3164 \\
IMF$_2$ & $5.14 \times 10^2$ & 0.1433 & 1075.0 & 0.2728 \\
IMF$_3$ & $1.19 \times 10^3$ & 0.3307 & 332.0 & 0.3769 \\
IMF$_4$ & $2.77 \times 10^3$ & 0.7724 & 151.0 & 0.5758 \\
IMF$_5$ & $6.19 \times 10^3$ & 1.7223 & 69.0 & 0.8446 \\
IMF$_6$ & $8.87 \times 10^3$ & 2.4622 & 34.0 & 1.0746 \\
IMF$_7$ & $3.82 \times 10^4$ & 10.5136 & 13.0 & 2.2015 \\
IMF$_8$ & $2.44 \times 10^5$ & 68.0636 & 4.0 & 6.0934 \\
\bottomrule
\end{tabular}
\end{table}

The analysis of the metrics reveals clear patterns in the structure of the IMFs. Lower-order IMFs (IMF$_1$ to IMF$_4$) exhibit high dominant frequencies (from 151 to 1287 cycles) and relatively small amplitudes (all below 0.58), indicating that they capture short-term oscillations and high-frequency noise. As the order of the IMF increases, a systematic decrease in dominant frequency and a progressive increase in variance and average amplitude are observed. Higher-order IMFs (IMF$_7$ and IMF$_8$) present dominant frequencies of 13 and 4 cycles and substantially larger average amplitudes (2.20 and 6.09), reflecting slower, larger-amplitude oscillations at the end of the intrinsic-mode ladder. Among the eight IMFs, IMF$_8$ attains the highest energy ($2.44 \times 10^5$) and variance (68.06), while most of the total signal energy is carried by the residue (not listed in Table~\ref{tab:imf_metrics}), consistent with interpreting the residue as the slow-moving price trend rather than as a fast oscillatory IMF.

\subsubsection{Validation of the decomposition}

Reconstruction was validated by comparing the sum of all IMFs and the residue with the original MSCI World closing-price series. In a consistent decomposition, the original MSCI Index and fully reconstructed signal must be identical.

The CEEMDAN residual is defined as the original signal minus the sum of all eight IMFs. For this sample it has mean 89.93, standard deviation 36.76, and range from 46.20 to 159.99. Figure~\ref{fig:imf_decomposition} (lower panel) contrasts this quantity with the analogous EEMD reconstruction gap on the same prices: the CEEMDAN trace is visually an almost noise-free, monotonically increasing trend, whereas EEMD retains a jagged slow component. That pattern is consistent with attributing the equity-like drift to a single smooth mode under CEEMDAN rather than folding fast oscillations into the baseline. The side-by-side comparison on the MSCI World series thus provides empirical evidence that CEEMDAN yields better decomposition outcomes than EEMD for this index, in the sense of a cleaner separation of slow drift from oscillatory structure on the same closing-price record. Because the CEEMDAN residual behaves here as an almost smooth, monotone trend, it can be modeled to a useful approximation by simple linear regression (for example, regressing the residual on a deterministic time or observation index) when a parsimonious parametric summary of the slow level is needed.

\begin{figure}[H]
    \centering
    \begin{tikzpicture}
        \node[anchor=south west,inner sep=0] (image) at (0,0) {\includegraphics[width=\textwidth]{./images/english/imf_decomposition.png}};
    \end{tikzpicture}
    \caption{MSCI World closing prices: CEEMDAN versus EEMD when the residue is excluded from the IMF sum. \textbf{Top:} original series and sum of the eight CEEMDAN IMFs (residue excluded). \textbf{Bottom:} $|x(t)-\sum_k \mathrm{IMF}_k(t)|$ without the residue for CEEMDAN ($K=8$) and for EEMD on the same series. The CEEMDAN trace is an almost smooth, monotonically increasing trend (magnitude of the CEEMDAN residue), whereas the EEMD trace is rougher and exhibits fine-scale undulations, indicating residual high-frequency content in the slow component under input-only ensemble averaging.}
    \label{fig:imf_decomposition}
\end{figure}

\FloatBarrier

\subsection{Graph transformation}

Graph transformations were performed via visibility and recurrence algorithms for all IMFs (and the residue) produced by the CEEMDAN decomposition above.

To visually illustrate the resulting structure of these transformations, Figure~\ref{fig:hvg_imf} shows the HVG graph obtained from IMF$_1$, which represents the highest frequency component of the decomposition. The visualization reveals the characteristic topology of HVG graphs: a sparse structure where each node (corresponding to a temporal point of the series) connects only with those nodes that satisfy the horizontal visibility criterion. This structure reflects the high-frequency nature of IMF$_1$, where rapid oscillations generate localized connectivity patterns that preserve the temporal information of the original series. The graphical representation allows us to appreciate how the visibility algorithm captures the geometric relationships between data points, transforming the one-dimensional temporal sequence into a two-dimensional network structure that maintains the structural properties of the modal component.

\begin{figure}[H]
    \centering
    \includegraphics[width=0.8\textwidth]{./images/english/hvg_imf.png}
    \caption{Horizontal visibility graph resulting from the transformation of IMF$_1$ of the MSCI World index. Each node represents a temporal point of the series, and the connections reflect the horizontal visibility criterion.}
    \label{fig:hvg_imf}
\end{figure}

\FloatBarrier

Specifically for the construction of recurrence graphs, optimal embedding parameters were determined independently for each IMF using the methods described in Section~\ref{sec:methodology}. The delay parameter $\tau$ was selected through mutual information analysis, identifying the first minimum of the mutual information function between the original series and its time-delayed version. The embedding dimension $d$ was estimated using the false nearest neighbors (FNN) method, systematically increasing the dimension until the proportion of false neighbors fell below the 10\% threshold. The recurrence threshold $\varepsilon$ was established as the 10th percentile of the distribution of pairwise Euclidean distances between all embedding vectors, ensuring that the graph captures the most significant recurrent structures. The parameter selection results reveal considerable variability among the different IMFs: several high-frequency IMFs select $d=4$, while progressively lower-frequency components tend to reduce the embedding dimension and increase the delay. Table~\ref{tab:parametros_recurrencia} presents the selected parameters for each IMF and the residue.

\begin{table}[h]
\centering
\caption{Embedding parameters and threshold $\varepsilon$ by IMF}
\label{tab:parametros_recurrencia}
\small
\begin{tabular}{lccc}
\toprule
\textbf{Component} & \textbf{$\tau$} & \textbf{$d$} & \textbf{$\varepsilon$} \\
\midrule
IMF$_1$ & 9 & 4 & 0.1847 \\
IMF$_2$ & 2 & 4 & 0.1561 \\
IMF$_3$ & 3 & 4 & 0.1348 \\
IMF$_4$ & 7 & 4 & 0.1887 \\
IMF$_5$ & 15 & 4 & 0.2562 \\
IMF$_6$ & 29 & 3 & 0.1891 \\
IMF$_7$ & 34 & 2 & 0.0798 \\
IMF$_8$ & 32 & 2 & 0.1216 \\
Residue & 23 & 2 & 0.1273 \\
\bottomrule
\end{tabular}
\end{table}

\subsubsection{General structural properties. Connectivity and distance}

The comparative analysis of the fundamental structural properties of graphs generated through the three transformation methods (NVG, HVG, and recurrence) reveals significant differences in the topology of the resulting networks. These differences reflect the inherent characteristics of each transformation algorithm and its interaction with the dynamic properties of the different IMFs.

In terms of network size, visibility-based graphs are built on the full daily grid (3587 nodes for the analyzed closing-price sample). Recurrence graphs are constructed on embedded state vectors, so the number of nodes is slightly smaller (roughly between 3529 and 3581 nodes in this study, depending on $(\tau,d)$). This difference is expected and does not prevent a meaningful comparison of qualitative structural patterns across methods.

The most distinctive characteristic among transformation methods is connection density, measured through the number of links and graph density. HVG graphs are extremely sparse, with densities between approximately $0.00057$ and $0.00111$, corresponding to between 3642 and 7161 undirected edges in the reported construction. Recurrence graphs have moderate densities between approximately $0.00255$ and $0.00287$, with between about 16188 and 18153 undirected edges under the 10th-percentile distance threshold used here. NVG graphs span the widest density range, from about $0.00166$ (high-frequency IMFs) up to about $0.766$ for the trend-dominated residue, with between about 10695 and 4927305 undirected edges depending on the IMF.

The global connectivity of graphs, measured through the number of connected components, reveals another fundamental difference among methods. Both HVG and NVG graphs are connected in this experiment (a single component), which facilitates message passing in many GNN pipelines. Recurrence graphs can fragment: the number of connected components ranges from 1 (notably for the residue in this run) up to about 1590, depending on the IMF. This fragmentation motivates strategies such as pooling over connected components or learned readouts that do not assume a single connected graph.

The average degree follows the same qualitative ordering: HVG graphs have low average degrees (about 2.0--4.0), recurrence graphs have roughly stable moderate average degrees (about 9.1--10.2 under the present thresholding), and NVG graphs vary widely (from single digits up to thousands for the densest modes), reflecting strong scale dependence in natural visibility.

In summary, the three transformations generate graph structures with distinctive topological properties. On one hand, HVG graphs are extremely sparse and always connected, while recurrence graphs present moderate densities but multiple disconnected components. Finally, NVG graphs show highly variable densities that depend strongly on IMF order. These structural differences have important implications for the design of GNN architectures and the selection of appropriate modeling strategies for each type of transformation. Table~\ref{tab:propiedades_estructurales} presents a summary of all structural properties obtained on the graphs after their transformation.

\begin{table}[h]
\centering
\caption{Summary of structural properties of transformed graphs}
\label{tab:propiedades_estructurales}
\small
\begin{tabular}{lccc}
\toprule
\textbf{Property} & \textbf{HVG} & \textbf{Recurrence} & \textbf{NVG} \\
\midrule
Density (range) & $0.00057$ - $0.00111$ & $0.00255$ - $0.00287$ & $0.00166$ - $0.766$ \\
Number of links (range) & $3642$ - $7161$ & $16188$ - $18153$ & $10695$ - $4927305$ \\
Connected components & $1$ (always) & $1$ - $1590$ & $1$ (always) \\
Diameter (range) & $24$ - $3558$ & $1725$ (single-component case) & $10$ - $242$ \\
Average degree (range) & $2.03$ - $3.99$ & $9.1$ - $10.2$ & $6.0$ - $2746$ \\
\bottomrule
\end{tabular}
\end{table}

\subsubsection{Clustering coefficients and local structure}

The clustering coefficient, which quantifies the tendency of nodes to form triangles or local cliques, addresses the question of how cohesive local neighborhoods are, revealing notable contrasts among transformation methods.

A direct comparison of the \emph{average} clustering coefficient ranges in Table~\ref{tab:clustering_coefficients} places NVG highest among the three (across IMF$_1$--IMF$_8$), followed by recurrence graphs, and finally HVG graphs---with the important caveat that NVG becomes extremely dense for the residue (cf.\ Table~\ref{tab:propiedades_estructurales}), which mechanically changes triangle counts, and the global clustering statistics for that NVG were not computed here because of cost.

However, this general hierarchy conceals subtler patterns that depend on IMF order under the CEEMDAN-based decomposition. For NVG graphs on IMF$_1$--IMF$_8$, the average local clustering coefficient ranges from about $0.483$ (IMF$_7$) to $0.754$ (IMF$_2$). The largest values occur at IMF$_2$ and IMF$_1$ ($\approx 0.752$--$0.754$), while IMF$_1$--IMF$_4$ span roughly $0.653$--$0.754$, indicating tight local neighborhoods at high frequency. From IMF$_5$ onward averages drift downward to a minimum near IMF$_7$, before rising again for IMF$_8$ ($\approx 0.607$), so the dependence on IMF index is not monotone. The residue NVG ($\approx 4.9 \times 10^6$ undirected edges) is omitted from the NVG ranges in Table~\ref{tab:clustering_coefficients}, and its regime is dominated by massive local connectivity.

The behavior of HVG graphs contrasts with NVG. IMF$_1$--IMF$_4$ retain comparatively high average clustering (about $0.529$--$0.658$) despite the horizontal-visibility restriction. From IMF$_5$ to IMF$_8$ averages decrease monotonically to about $0.389$. The residue HVG drops to about $0.008$, i.e.\ negligible triangle participation in the mean, consistent with a near-chain visibility pattern on a smooth, slowly varying component---not a complete absence of local clustering at every node, but a strong shift in the distribution. Standard deviations of the node-level clustering coefficients span about $0.065$--$0.289$, largest for the highest-frequency IMFs. Medians range from $0$ (residue) to about $0.667$ (IMF$_1$--IMF$_2$), with several intermediate IMFs at $0.5$, indicating a mixture of nodes with zero local clustering and nodes with moderate clustering.

Recurrence graphs do not follow the same IMF-order trend as HVG. The lowest average clustering appears for IMF$_3$ and IMF$_4$ (about $0.211$--$0.221$), while IMF$_6$, IMF$_8$, and the residue reach the upper part of the range (about $0.409$--$0.459$). This pattern suggests that, under the chosen recurrence construction, mid-frequency embeddings can be the least triangle-rich on average, whereas slower oscillatory modes and the residue align with higher mean local clustering. Standard deviations of the per-node coefficients lie between about $0.254$ and $0.326$. Medians vary widely: some IMFs exhibit a median of $0$ (many embedded states carry zero local clustering), while others reach medians up to about $0.639$.

Table~\ref{tab:clustering_coefficients} synthesizes these results, showing the ranges of average, median, and standard deviation of the \emph{node-level} clustering coefficients for each method (NVG ranges use IMF$_1$--IMF$_8$ only, with the residue NVG omitted).

\begin{table}[h]
\centering
\caption{Clustering coefficients by transformation method and IMF order}
\label{tab:clustering_coefficients}
\small
\begin{tabular}{lccc}
\toprule
\textbf{Method} & \textbf{Average coefficient} & \textbf{Median coefficient} & \textbf{Standard deviation} \\
 & \textbf{(range)} & \textbf{(range)} & \textbf{(range)} \\
\midrule
HVG & $0.0083$ - $0.6584$ & $0.0000$ - $0.6667$ & $0.0652$ - $0.2895$ \\
Recurrence & $0.2109$ - $0.4585$ & $0.0000$ - $0.6389$ & $0.2536$ - $0.3262$ \\
NVG & $0.4832$ - $0.7536$ & $0.4991$ - $0.8333$ & $0.1937$ - $0.3074$ \\
\bottomrule
\end{tabular}
\end{table}

The local structure of transformed graphs, characterized through clustering coefficients, reveals distinctive topological properties that have important implications for GNN modeling. NVG graphs, with high average clustering on oscillatory IMFs (but an omitted residue NVG here), present locally cohesive neighborhoods where message passing can emphasize short-range structure. HVG graphs combine strong clustering at high frequency with a sharp collapse toward near-zero mean clustering on the residue. Recurrence graphs yield intermediate averages with a non-monotone IMF dependence, so modeling choices should reflect embedding dimension, thresholding, and possible fragmentation (as in the connectivity discussion above) alongside these clustering summaries.

\subsubsection{Centrality analysis}

Identifying which nodes play critical roles in the graph structure is fundamental for understanding how information flows and where influence is concentrated. For this purpose, we examine three centrality metrics that capture complementary aspects: betweenness centrality identifies strategic bridges, closeness centrality reveals nodes with rapid access to the rest of the network, and eigenvector centrality measures influence propagated through connections. Unless stated otherwise, summaries refer to the mean of the per-node centralities (NetworkX definitions, normalized betweenness). For recurrence graphs with many edges, betweenness and closeness use the library's sampling-based estimates. For dense NVGs, betweenness uses random $k$-node sampling and closeness uses a random subset of source nodes, both with fixed seeds in the reproducibility script. NVG edge lists for IMF$_1$--IMF$_8$ are loaded from a prior export to avoid rebuilding visibility graphs when possible. The residue NVG is omitted from centrality summaries here because it exceeds the edge budget used in the pipeline.

When analyzing which nodes act as critical bridges, contrasts emerge among transformation methods. \textbf{Betweenness} means stay low on NVGs over IMF$_1$--IMF$_8$ (about $0.0010$--$0.0198$), consistent with dense graphs that spread shortest-path traffic across many nodes. Node-level maxima can still reach about $0.60$ on the highest-frequency NVGs. HVGs show a much wider mean range (about $0.0028$--$0.333$). IMF$_1$ remains near $0.0028$, whereas the residue reaches about $0.333$, reflecting a near-path visibility structure where a handful of dates can act as bottlenecks. Maximum per-node betweenness on HVGs peaks near $0.63$--$0.65$ on several IMFs (including IMF$_1$ and IMF$_8$), far above typical NVG maxima on comparable scales. Recurrence graphs span about $0.0003$--$0.173$ in mean betweenness. The smallest means occur for mid-frequency embeddings (e.g.\ IMF$_3$), while the residue attains the largest mean in this run, consistent with a single large component where a few regions still carry disproportionate inter-cluster traffic under approximate estimation.

\textbf{Closeness} means are highest for NVGs on IMF$_1$--IMF$_8$ (about $0.020$--$0.21$ in this computation, with IMF$_4$ near the upper end and the most oscillatory IMF$_8$ NVG near the lower end), reflecting short average distances in dense visibility graphs. HVG closeness means range from about $0.0009$ on the residue to about $0.091$ on IMF$_1$, while recurrence graphs span about $0.0017$--$0.082$, with IMF$_2$ toward the top and the residue toward the bottom. These patterns align with fragmentation and long effective distances in some recurrence constructions.

\textbf{Eigenvector} centrality means remain modest for HVGs (about $0.0013$--$0.0059$) and recurrence graphs (about $0.0029$--$0.011$), whereas NVG means lie slightly higher (about $0.0051$--$0.011$). Maxima tell a complementary story. HVG eigenvector peaks near $0.62$ on IMF$_8$ despite a low mean, highlighting a few locally influential dates on an otherwise sparse graph. NVG maxima are smaller in relative terms (often below $0.40$ on IMF$_1$--IMF$_3$), while recurrence maxima stay below about $0.15$ except where a dominant component concentrates weight.

\begin{table}[h]
\centering
\caption{Centrality measures by transformation method}
\label{tab:centralidad}
\small
\begin{tabular}{lccc}
\toprule
\textbf{Metric} & \textbf{HVG} & \textbf{Recurrence} & \textbf{NVG} \\
 & \textbf{(average range)} & \textbf{(average range)} & \textbf{(average range)} \\
\midrule
Betweenness & $0.0028$ - $0.3332$ & $0.0003$ - $0.1732$ & $0.0010$ - $0.0198$ \\
Closeness & $0.0009$ - $0.0914$ & $0.0017$ - $0.0815$ & $0.0203$ - $0.2085$ \\
Eigenvector & $0.0013$ - $0.0059$ & $0.0029$ - $0.0107$ & $0.0051$ - $0.0113$ \\
\bottomrule
\end{tabular}
\end{table}

Table~\ref{tab:centralidad} synthesizes the observed ranges of \emph{mean} per-node centralities for each method (NVG column: IMF$_1$--IMF$_8$ only, with the residue NVG omitted). These differences matter for GNN design. NVGs support diffuse, short-range summaries. HVGs can concentrate bridging on a narrow temporal backbone, especially on the residue. Recurrence graphs mix low global closeness with occasional high mean betweenness on the residue, so readouts should tolerate disconnectedness and approximation noise in the chosen estimators.

\subsubsection{Patterns according to IMF order}

Under the CEEMDAN decomposition used here (IMF$_1$--IMF$_8$ plus the residue), IMF index orders components from faster oscillations toward the slow trend. Even so, the graph statistics reviewed above do not track that index in a simple monotone way. Each construction folds scale-dependent dynamics into a different topology, which means ``high-frequency'' and ``low-frequency'' are helpful labels, yet they cannot substitute for component-wise inspection.

Consider first the \textbf{HVG} graphs. From IMF$_1$ toward IMF$_8$, diameters tend to grow and mean clustering drifts downward (see the clustering discussion and Table~\ref{tab:clustering_coefficients}), whereas connectivity remains trivial in the sense of a single component throughout (Table~\ref{tab:propiedades_estructurales}). Against that gradual trend, the residue stands apart: diameter explodes, mean clustering is essentially negligible, and mean betweenness rises sharply compared with the early IMFs (Table~\ref{tab:centralidad}), which is what one expects when horizontal visibility sees a smooth series as an almost one-dimensional chain with a handful of bottleneck dates. For practitioners, the upshot is a clear shift from fairly compact, triangle-rich HVGs at high frequency to a sparse, elongated graph on the residue where bridging can concentrate on very few nodes.

\textbf{Recurrence graphs} paint a contrasting picture. Here, fragmentation is tightly tied to the IMF-specific triple $(\tau,d,\varepsilon)$ in Table~\ref{tab:parametros_recurrencia}. Many modes split into hundreds or more than a thousand components, yet the residue may still appear connected under the same recipe (Table~\ref{tab:propiedades_estructurales}). Nor does clustering follow IMF order in lockstep. Mid-frequency embeddings can be the least triangle-rich on average, while slower oscillations or the residue may sit higher (clustering discussion above). Centrality behaves similarly. Mean betweenness is modest on several mid-frequency IMFs but peaks on the residue in this run (Table~\ref{tab:centralidad}). In practice, it is safer to treat recurrence graphs as a \emph{family} of regimes than as one uniformly ``cohesive'' class. Whenever components multiply, encoders or pooling that work per piece are natural, and any summary that presumes a single connected graph should be checked mode by mode.

Turning finally to \textbf{natural visibility graphs (NVG)}, they exhibit the widest cross-IMF spread in density, degree, diameter, and clustering (Tables~\ref{tab:propiedades_estructurales}--\ref{tab:clustering_coefficients}). IMF$_1$--IMF$_4$ pair comparatively moderate edge counts with strong mean clustering, while deeper IMFs, including IMF$_8$, sit in an intermediate regime (still dense, but far from the residue). At the bottom of the scale, the residue NVG is orders of magnitude heavier than any IMF (Table~\ref{tab:propiedades_estructurales}), to the point that several global summaries were skipped earlier for computational reasons. That limitation is part of the modeling problem itself whenever one insists on a full NVG at the lowest frequency.

In summary, IMF order is best read as a guide to \emph{interpretation}, not as a universal knob. HVG compactness tends to erode toward the residue. Recurrence links embedding choices to fragmentation and to local structure that can rise or fall in ways that are not monotone in IMF index. NVG couples scale to explosive variation in density. For graph-based learning, a workable rule of thumb is that \emph{early IMFs} often pair well with straightforward message passing on connected HVG/NVG views where neighborhoods still carry interpretable structure. \emph{Mid-frequency IMFs} are frequently where recurrence fragmentation and non-monotone clustering force explicit choices about components and thresholds. \emph{The residue} is where path-like HVGs, possibly dominant recurrence betweenness when the graph is connected, and prohibitively dense NVGs all meet, so sparsification, sampling, or multi-scale pooling stop being optional refinements and become part of the design.

\subsubsection{Implications for GNN modeling}

The structural properties identified in transformed graphs have important implications for the application of graph neural networks. NVG graphs, with their high density and very short-distance structure, provide a rich environment for message passing in GNNs, allowing information to propagate efficiently through the network. However, high density can also increase computational complexity and potentially introduce noise in learned representations.

HVG graphs behave differently: their sparser structure yields a leaner representation that reduces computational overhead. Still, the wide range of distances observed—particularly in low-frequency IMFs—can complicate the learning of stable representations. One clear benefit is that all HVG graphs remain connected, enabling global message passing without resorting to specialized techniques for handling isolated components.

Recurrence graphs introduce another layer of complexity. Most IMFs contain numerous disconnected components, forcing GNN designers to adopt component-wise processing or pooling methods that merge information across components. Despite this challenge, recurrence graphs encode core dynamic features of the underlying system, making them useful for prediction tasks.

Taken together, these three methods complement each other. A hybrid approach that integrates multiple graph representations could exploit the advantages of each while compensating for their weaknesses. This direction opens up promising avenues for developing multi-view GNN architectures tailored to financial time series decomposed through CEEMDAN.

\FloatBarrier

\subsection{Dimensionality reduction of the oscillatory IMF panel}
\label{sec:dim_red_imf}

The graph analyses above use the full eight oscillatory IMF channels (plus the residue) on the daily grid, which for some constructions implies a large number of nodes and, especially for NVGs at lower frequency, a very large edge budget. One therefore asks whether those eight channels can be replaced by a smaller set of derived coordinates \emph{at each time index} so that downstream graph construction or repeated experiments could operate on a compact multichannel state. The linear PCA and ICA pipelines used here are specified in Section~\ref{sec:method_dim_red}. This subsection reports only the empirical application to the MSCI CEEMDAN panel.

The same workflow was also run for ranks $k\in\{3,5\}$. Figures and tabulated reconstruction metrics below foreground $k=4$ as a representative mid-rank compromise between compression and fit.

For the MSCI World CEEMDAN panel aligned with the empirical IMF table ($T=3596$ daily observations), the eight oscillatory IMF columns were standardized columnwise and projected to $k=4$ scores using FastICA and, separately, PCA. The CEEMDAN residue was excluded from this linear block so that the slow trend is not folded into the reduced coordinates. Figure~\ref{fig:imf_orig_vs_red} shows the four ICA score series $Z_1$--$Z_4$ over the sample as an illustration of the compact multichannel state. PCA score trajectories can differ at the same rank even when the induced reconstruction agrees with ICA (see below).

\begin{figure}[H]
    \centering
    \includegraphics[width=0.72\textwidth]{./images/english/imf_original_vs_reduced_fastica_k4.png}
    \caption{MSCI World CEEMDAN oscillatory block ($T=3596$): four reduced coordinates $Z_1$--$Z_4$ from ICA after column standardization of IMF$_1$--IMF$_8$ ($k=4$). The CEEMDAN residue is excluded from this linear reduction.}
    \label{fig:imf_orig_vs_red}
\end{figure}

Scores were mapped back through the inverse linear operator in score space and the inverse column standardization so that the fitted oscillatory block $\hat{\mathbf{X}}$ is expressed in the original IMF measurement units. Working in that native domain is deliberate: reconstruction error and information loss after reduction are read by comparing $\hat{\mathbf{X}}$ to $\mathbf{X}$ mode by mode with RMSE and $R^2$ on the CEEMDAN scale rather than only in abstract standardized coordinates. Figure~\ref{fig:imf_orig_vs_recon_panel} plots each $\mathrm{IMF}_j$ alongside $\widehat{\mathrm{IMF}}_j$ from that inverse map. For this sample and rank, ICA and PCA yield \emph{numerically identical} reconstructions $\hat{\mathbf{X}}$, i.e.\ the same rank-four orthogonal projector of the standardized block, while the four score series $Z_j$ need not coincide between methods. In aggregate, $\|\mathbf{X}-\hat{\mathbf{X}}\|_F/\|\mathbf{X}\|_F\approx 0.59$ and the arithmetic mean of the eight columnwise coefficients $R^2_j$ between $\mathrm{IMF}_j$ and $\widehat{\mathrm{IMF}}_j$ is about $0.64$, consistent with the variance captured by the leading four principal directions after column scaling.

\begin{figure}[H]
    \centering
    \includegraphics[width=\textwidth]{./images/english/imf_original_vs_reconstructed_k4.png}
    \caption{MSCI World ($T=3596$). \textbf{Left:} CEEMDAN IMF$_1$--IMF$_8$. \textbf{Right:} rank-$k=4$ linear reconstructions $\widehat{\mathrm{IMF}}_j$ from the standardized oscillatory block (residue excluded).}
    \label{fig:imf_orig_vs_recon_panel}
\end{figure}

Table~\ref{tab:imf_red_recon_k4} reports RMSE and $R^2$ between each native mode and its reconstruction.

\begin{table}[h]
\centering
\caption{Rank-$k=4$ linear reconstruction of each oscillatory IMF on the MSCI sample ($T=3596$): RMSE and $R^2$ between $\mathrm{IMF}_j$ and $\widehat{\mathrm{IMF}}_j$ (ICA and PCA coincide numerically for $\widehat{\mathrm{IMF}}_j$).}
\label{tab:imf_red_recon_k4}
\small
\begin{tabular}{lcc}
\toprule
\textbf{Mode} & \textbf{RMSE} & \textbf{$R^2$} \\
\midrule
IMF$_1$ & 0.380 & 0.300 \\
IMF$_2$ & 0.227 & 0.639 \\
IMF$_3$ & 0.319 & 0.663 \\
IMF$_4$ & 0.488 & 0.747 \\
IMF$_5$ & 0.678 & 0.724 \\
IMF$_6$ & 0.811 & 0.724 \\
IMF$_7$ & 1.789 & 0.642 \\
IMF$_8$ & 4.876 & 0.650 \\
\bottomrule
\end{tabular}
\end{table}

IMF$_1$ shows the weakest $R^2$ (about $0.30$, with RMSE $\approx 0.38$), so most of its variance lies outside the retained four-dimensional subspace. Mid-frequency modes reach $R^2$ up to roughly $0.75$ (IMF$_4$) with RMSE below $0.7$ through IMF$_6$. For IMF$_7$ and IMF$_8$, $R^2$ remains near $0.64$--$0.65$ but RMSE rises sharply ($\approx 1.8$ and $\approx 4.9$), reflecting wider amplitude on the slowest oscillations. Here absolute pointwise errors can be large while the linear fit stays only moderate on a relative scale.

Section~\ref{sec:synthetic_emd} already measures cross-IMF linear coupling through $\bar{\rho}$ and $\pi_{0.05}$ (Table~\ref{tab:synthetic_emd_summary}), partly to justify CEEMDAN as yielding comparatively modest average Pearson redundancy among native modes on the calibrated stress test. Compressing eight IMF channels to four scores $Z_j$ breaks the one-to-one alignment with mode indices, so it is natural to ask whether the $Z_j$ inherit a similar qualitative property: limited mutual linear dependence across parallel channels rather than a few scores that move together like a low-rank copy of the same component. Let $r_{ij}$ denote the sample Pearson correlation between $Z_i$ and $Z_j$ over the $T$ time points, for $1\le i<j\le 4$. For the standardized scores from the PCA and ICA fits above, linear decorrelation of the $Z_j$ is the design target. Empirically, all $|r_{ij}|$ are below $10^{-12}$ (numerical roundoff only). Using the same uncorrected Pearson test protocol as in Table~\ref{tab:synthetic_emd_summary}, none of the six pairs yields $p<0.05$, so $\pi_{0.05}=0$ for this diagnostic. Spearman correlations among the $Z_j$ are not constrained to vanish when Pearson correlations do. Their absolute values average about $0.035$ and reach a maximum near $0.07$, indicating only weak monotonic co-movement while linear decorrelation remains essentially perfect. Higher-order dependence is therefore only partially controlled relative to the ideal ICA objective.

\FloatBarrier

\section{Conclusions and future lines}
\label{conclfutu}

\subsection{Conclusions}
As a result of this study, a methodological framework has been developed that combines EMD with graph transformation techniques to analyze financial time series, specifically applied to the MSCI World index. The results obtained show that this integration of complementary techniques enables a deeper understanding of the complex, nonlinear, and non-stationary dynamics that characterize financial markets.

Statistical validation using the ADF, KPSS, and BDS tests confirmed that the MSCI World index satisfies the necessary conditions to apply EMD effectively, establishing a solid foundation for subsequent analysis. A controlled comparison of EMD, EEMD, and CEEMDAN on synthetic signals (Section~\ref{sec:synthetic_emd}), including a multi-component benchmark with an explicit upward drift and ensemble hyperparameters tuned on that benchmark, illustrated separation limits, mode mixing, and the practical differences between input-only noise injection and stage-wise complete ensemble schemes, motivating the use of CEEMDAN on empirical data. Decomposition using CEEMDAN produced eight IMFs plus a residue, each capturing different temporal scales of market dynamics. Figure~\ref{fig:imf_decomposition} contrasts CEEMDAN with EEMD on the same prices and shows that, when oscillatory IMFs are summed without the residue, CEEMDAN isolates a smooth monotone trend in the implied gap, whereas EEMD leaves noticeable short-term ripple in that quantity, supporting the choice of CEEMDAN for the graph stage. Most energy is concentrated in the residue (slow trend of the price level), while the eight IMFs describe the oscillatory structure around that trend at increasing scales. A complementary linear reduction (Section~\ref{sec:method_dim_red} and empirical results in Section~\ref{sec:dim_red_imf}), positioned after the graph discussion to link compression with large node and edge regimes, illustrated ICA- and PCA-based compression of the eight oscillatory IMF channels to $k\in\{3,4,5\}$ scores per time index, with per-mode reconstruction error and cross-score dependence assessed analogously to the synthetic IMF diagnostics. The empirical graph stage used the native IMF panel without applying that reduction.

The transformation of these IMFs into graph representations through three complementary methods reveals distinctive topological properties that reflect different aspects of the underlying structure of financial data. These properties can be summarized as follows:
\begin{itemize}
    \item HVG graphs generate extremely sparse structures with low density and average degree while maintaining full connectivity. However, they exhibit high variability in distances for low-frequency IMFs, which complicates the learning of consistent representations. 
    \item Recurrence networks show moderate edge densities in the connected regime, but often split into many disconnected components, with counts varying widely across IMFs under the chosen threshold. This fragmentation requires architecture choices (e.g.\ per-component encoders or pooling) that do not assume a single connected graph. 
    \item NVG graphs exhibit the greatest variability in density, ranging from sparse structures for high-frequency IMFs to extremely dense graphs for low-frequency IMFs. This high density, while capturing rich relationships, significantly increases computational complexity.
\end{itemize}


A key finding is that graph statistics relate to IMF scale, yet they need not track IMF index monotonically within a given transformation. HVGs tend toward degraded compactness and very low clustering on the residue. Recurrence graphs couple embedding parameters to fragmentation and IMF-dependent local structure. NVGs exhibit strong scale-dependent density, with an extremely heavy residue graph. This diversity suggests that no single graph representation is uniformly optimal across IMF components.

The results of this study provide crucial insights for the design of GNN architectures applied to financial time series forecasting. The identified topological properties (extreme sparsity in HVGs, disconnected components in recurrence networks, and high density in NVGs) pose specific modeling challenges that must be addressed through specialized architectural strategies. In particular, hybrid approaches are required that can leverage the complementary strengths of multiple graph representations simultaneously.

In conclusion, this work establishes that the integration of EMD with graph analysis constitutes a powerful tool for unraveling the multilevel structure of financial time series. The discovered topological patterns not only deepen our understanding of market dynamics, but also inform the development of more sophisticated GNN-based forecasting systems. The proposed methodology is generalizable to other financial indices and asset classes, opening promising avenues for future research in quantitative financial market analysis.

\subsection{Future lines}
Limitations surfaced here double as a roadmap. One priority is to hard-wire GNN architectures around the topological signatures observed for each IMF graph family. NVG, HVG, and recurrence stress different geometry, so multi-view stacks that ingest several graph views jointly \citep{liu2023multimodal, wang2023hatr} look especially relevant. Adaptive attention could then reweight views as the forecast horizon or volatility regime shifts \citep{su2024attention}.

Recent work already couples temporal encoders with graph convolutions to track both intra-series evolution and cross-asset structure \citep{liu2023multimodal, wang2023hatr, jin2023temporal}. IMF panels add another axis: edges should arguably move with regimes rather than stay static. Recurrence plots here hint at higher-order dependence across modes, so hypergraph constructions \citep{su2024attention, agarwal2022think} deserve prototypes beyond pairwise edges. Concrete architectural bets worth testing next include hyperbolic hypergraph attention for multiscale IMF hierarchies \citep{agarwal2022think}, heterogeneous graphs that type edges by relation \citep{jin2023temporal}, and dual attention over nodes versus relation types to tame mixed signals \citep{wang2023hatr}.

Scalability matters once the universe of assets and refresh cadence grows. Distributed graph engines with partitioning, multilevel parallelism, and communication-aware schedulers already underpin large GNN training \citep{xu2025capturing}. Translating those recipes to CEEMDAN-driven IMF graphs will require partitions that honor temporal adjacency, bounded halo exchange between workers, batched windows for data parallelism, and caches for quasi-stable IMFs that change slowly.

Streams arrive continuously, so graphs built on rolling CEEMDAN outputs need maintenance policies that combine incremental sifting updates instead of full recomputes whenever prices tick, dynamic-graph GNN layers that accept edge insertions and deletions \citep{gravina2023scalable, liu2023decoupled}, and latency-aware serving stacks for live desks. Dense NVGs and fragmented recurrence graphs also stress memory, and sparsifying transforms, neighborhood-importance sampling, and asynchronous training collectives belong on the same roadmap. When densification starts before message passing, widening the pilot in Section~\ref{sec:method_dim_red} beyond linear PCA and FastICA is worth exploring through nonlinear autoencoders, manifold embeddings, sparse factorizations, or kernel lifts that could shrink the channel count before visibility or recurrence rules fire, trading reconstruction bias for thinner adjacency lists under fixed constructors.

Multi-index portfolios compound every scaling issue. Federated workflows could train shared encoders without centralizing order flows. Transfer models fit on liquid benchmarks might warm-start regional books. Aggregation layers could couple macro indices with sector peers instead of treating MSCI World in isolation.

Other research directions include:
\begin{itemize}
    \item Port the methodological framework to additional instruments and check whether the topological motifs generalize,
    \item Track how IMF graph statistics shift across market states and whether those shifts anticipate regime changes,
    \item Build post-hoc or intrinsic interpretability tools that attribute predictions to specific IMFs and graph motifs,
    \item Automate EEMD and CEEMDAN hyperparameter choice (ensemble size, noise scales) per series, extending the grid-search discipline used on the synthetic stress test,
    \item Compare further time-series-to-graph maps and fold in non-price channels (news, social sentiment) where licensing permits,
    \item Push IMF-panel reduction past the PCA and FastICA baseline in Section~\ref{sec:method_dim_red}: benchmark nonlinear and sparse families jointly with the graph constructor when the goal is to cap node and edge growth on each reduced channel.
\end{itemize}

Together, these threads span specialized graph encoders, systems for streaming and distributed training, and richer data interfaces. Crossing them with adaptive decomposition remains a fertile niche for financial forecasting and for risk workflows that must respect latency, privacy, and reconstruction fidelity at once.

\FloatBarrier

\begin{thebibliography}{99}

\bibitem{bollerslev1986generalized}
Bollerslev, T. (1986).
\newblock Generalized autoregressive conditional heteroskedasticity.
\newblock {\em Journal of Econometrics}, 31(3), 307--327.

\bibitem{brock1996test}
Brock, W.~A., Dechert, W.~D., Scheinkman, J.~A., \& LeBaron, B. (1996).
\newblock A test for independence based on the correlation dimension.
\newblock {\em Econometric Reviews}, 15(3), 197--235.

\bibitem{chen2025identifying}
Zeng, Z., \& Chen, Y. (2025).
\newblock Identifying and forecasting recurrently emerging stock trend structures via rising visibility graphs.
\newblock {\em Forecasting}, 7(2), 26.
\newblock \url{https://doi.org/10.3390/forecast7020026}

\bibitem{dickey1979distribution}
Dickey, D.~A., \& Fuller, W.~A. (1979).
\newblock Distribution of the estimators for autoregressive time series with a unit root.
\newblock {\em Journal of the American Statistical Association}, 74(366a), 427--431.

\bibitem{eckmann1987recurrence}
Eckmann, J.~P., Kamphorst, S.~O., \& Ruelle, D. (1987).
\newblock Recurrence plots of dynamical systems.
\newblock {\em EPL (Europhysics Letters)}, 4(9), 973.

\bibitem{engle1982autoregressive}
Engle, R.~F. (1982).
\newblock Autoregressive conditional heteroscedasticity with estimates of the variance of United Kingdom inflation.
\newblock {\em Econometrica}, 50(4), 987--1007.

\bibitem{han2022multiscale}
Han, M., Fan, Q., \& Ling, G. (2022).
\newblock Multiscale online-horizontal-visibility-graph correlation analysis of financial market.
\newblock {\em Physica A: Statistical Mechanics and Its Applications}, 607, 128195.
\newblock \url{https://doi.org/10.1016/j.physa.2022.128195}

\bibitem{hyvarinen2000independent}
Hyvärinen, A., \& Oja, E. (2000).
\newblock Independent component analysis: algorithms and applications.
\newblock {\em Neural Networks}, 13(4--5), 411--430.
\newblock \url{https://doi.org/10.1016/S0893-6080(00)00026-5}

\bibitem{jolliffe2016principal}
Jolliffe, I.~T., \& Cadima, J. (2016).
\newblock Principal component analysis: a review and recent developments.
\newblock {\em Philosophical Transactions of the Royal Society A}, 374(2065), 20150202.
\newblock \url{https://doi.org/10.1098/rsta.2015.0202}

\bibitem{huang1998empirical}
Huang, N.~E., Shen, Z., Long, S.~R., Wu, M.~C., Shih, H.~H., Zheng, Q., \ldots \& Liu, H.~H. (1998).
\newblock The empirical mode decomposition and the Hilbert spectrum for nonlinear and non-stationary time series analysis.
\newblock {\em Proceedings of the Royal Society of London. Series A: Mathematical, Physical and Engineering Sciences}, 454(1971), 903--995.

\bibitem{kwiatkowski1992testing}
Kwiatkowski, D., Phillips, P.~C., Schmidt, P., \& Shin, Y. (1992).
\newblock Testing the null hypothesis of stationarity against the alternative of a unit root: How sure are we that economic time series have a unit root?
\newblock {\em Journal of Econometrics}, 54(1-3), 159--178.

\bibitem{lacasa2008visibility}
Lacasa, L., Luque, B., Ballesteros, F., Luque, J., \& Nuño, J.~C. (2008).
\newblock From time series to complex networks: The visibility graph.
\newblock {\em Proceedings of the National Academy of Sciences}, 105(13), 4972--4975.

\bibitem{ljung1978measure}
Ljung, G.~M., \& Box, G.~E. (1978).
\newblock On a measure of lack of fit in time series models.
\newblock {\em Biometrika}, 65(2), 297--303.

\bibitem{rilling2008one}
Rilling, G., \& Flandrin, P. (2008).
\newblock One or two frequencies? The empirical mode decomposition answers.
\newblock {\em IEEE Transactions on Signal Processing}, 56(1), 85--95.

\bibitem{luque2009horizontal}
Luque, B., Lacasa, L., Ballesteros, F., \& Luque, J. (2009).
\newblock Horizontal visibility graphs: Exact results for random time series.
\newblock {\em Physical Review E}, 80(4), 046103.

\bibitem{song2024early}
Song, S., \& Li, H. (2024).
\newblock Early warning signals for stock market crashes: empirical and analytical insights utilizing nonlinear methods.
\newblock {\em EPJ Data Science}, 13(1), 16.
\newblock \url{https://doi.org/10.1140/epjds/s13688-024-00457-2}

\bibitem{su2025hybrid}
Su, J., Li, H., Wang, R., Guo, W., Hao, Y., Kurths, J., \& Gao, Z. (2025).
\newblock A hybrid dual-branch model with recurrence plots and transposed transformer for stock trend prediction.
\newblock {\em Chaos: An Interdisciplinary Journal of Nonlinear Science}, 35(1), 013125.
\newblock \url{https://doi.org/10.1063/5.0233275}

\bibitem{wu2009ensemble}
Wu, Z., \& Huang, N.~E. (2009).
\newblock Ensemble empirical mode decomposition: A noise-assisted data analysis method.
\newblock {\em Advances in Adaptive Data Analysis}, 1(1), 1--41.

\bibitem{torres2011complete}
Torres, M.~E., Colominas, M.~A., Schlotthauer, G., \& Flandrin, P. (2011).
\newblock A complete ensemble empirical mode decomposition with adaptive noise.
\newblock {\em 2011 IEEE International Conference on Acoustics, Speech and Signal Processing (ICASSP)}, 4144--4147.

\bibitem{liu2023multimodal}
Liu, R., Liu, H., Huang, H., Yang, J., Gao, Y., Zhu, J., Wang, C., \& Xu, K. (2023).
\newblock Multimodal multiscale dynamic graph convolution networks for stock price prediction.
\newblock {\em Pattern Recognition}, 144, 110211.
\newblock \url{https://doi.org/10.1016/j.patcog.2023.110211}

\bibitem{wang2023hatr}
Wang, H., Wang, T., Li, S., Li, J., \& Zhang, Z. (2023).
\newblock HATR-I: Hierarchical Adaptive Temporal Relational Interaction for Stock Trend Prediction.
\newblock {\em IEEE Transactions on Knowledge and Data Engineering}, 35(7), 7492--7506.
\newblock \url{https://doi.org/10.1109/TKDE.2022.3188320}

\bibitem{su2024attention}
Su, H., Wang, X., Qin, Y., Zhao, X., \& Lin, C. (2024).
\newblock Attention based adaptive spatial--temporal hypergraph convolutional networks for stock price trend prediction.
\newblock {\em Expert Systems with Applications}, 237, 121899.
\newblock \url{https://doi.org/10.1016/j.eswa.2023.121899}

\bibitem{jin2023temporal}
Jin, M., Koh, H.~Y., Wen, Q., Zambon, D., Alippi, C., Webb, G.~I., King, I., \& Pan, S. (2023).
\newblock Temporal and Heterogeneous Graph Neural Network for Financial Time Series Prediction.
\newblock In {\em Proceedings of the 31st ACM International Conference on Information and Knowledge Management} (CIKM '22), 3584--3593.
\newblock \url{https://doi.org/10.1145/3511808.3557089}

\bibitem{agarwal2022think}
Agarwal, S., Sawhney, R., Thakkar, M., \& Gupta, D. (2022).
\newblock THINK: Temporal Hypergraph Hyperbolic Network.
\newblock In {\em Proceedings of the 2022 IEEE International Conference on Data Mining} (ICDM), 1--10.
\newblock \url{https://doi.org/10.1109/ICDM54844.2022.00096}

\bibitem{gravina2023scalable}
Gravina, A., Zambon, D., Bacciu, D., \& Alippi, C. (2023).
\newblock Scalable spatiotemporal graph neural networks.
\newblock In {\em Proceedings of the AAAI Conference on Artificial Intelligence}, 37(6), 7688--7696.

\bibitem{liu2023decoupled}
Liu, Y., Li, K., Xie, Y., Guo, Y., Yan, J., \& Yu, P.~S. (2023).
\newblock Decoupled Graph Neural Networks for Large Dynamic Graphs.
\newblock {\em arXiv preprint arXiv:2305.08273}.
\newblock \url{https://doi.org/10.48550/arXiv.2305.08273}

\bibitem{xu2025capturing}
Xu, Q. (2025).
\newblock Capturing Structural Evolution in Financial Markets with Graph Neural Time Series Models.
\newblock {\em Preprints}, 2025071260.
\newblock \url{https://doi.org/10.20944/preprints202507.1260.v1}

\end{thebibliography}

\end{document}
